\documentclass[12pt]{article}
\usepackage{amsmath,amssymb,amsthm,geometry,booktabs,float,longtable,array}
\usepackage[hidelinks]{hyperref}
\usepackage[table]{xcolor}
\geometry{margin=1in}

\newtheorem{theorem}{Theorem}[section]
\newtheorem{lemma}[theorem]{Lemma}
\newtheorem{proposition}[theorem]{Proposition}
\newtheorem{corollary}[theorem]{Corollary}
\newtheorem{definition}[theorem]{Definition}
\newtheorem{remark}[theorem]{Remark}
\newtheorem{conjecture}[theorem]{Conjecture}

% Custom QED symbol: filled square with white QED punched through.
% At normal size it reads as a Halmos tombstone; zoom in to discover the letters.
\renewcommand{\qedsymbol}{%
  \raisebox{0.05ex}{%
    \rlap{\rule{0.7em}{0.7em}}%
    \hbox to 0.7em{\hfil%
      \raisebox{0.12em}{%
        \color{white}\fontsize{4.5}{4.5}\selectfont\bfseries\sffamily%
        \rlap{Q}\kern0.01em%
        \rlap{E}\kern0.04em%
        D%
      }%
    \hfil}%
  }%
}

\title{Two Foci, One Cup: Successive World-Record Equations for Ellipse Perimeter}
\author{Monir Mamoun}
\date{March 2026}

\begin{document}
\maketitle

%==========================================================================
\begin{abstract}
%==========================================================================

An ellipse has two foci, and they do not see eye to eye.

\medskip
\noindent\textbf{The problem.}
No finite algebraic expression gives the exact perimeter of an ellipse.  The best known approximation prior to this work is Ramanujan's second formula (1914), which we call \emph{Ramanujan~II} (abbreviated \textbf{RII} in Roman numerals, to distinguish it from the tower notation $R_k$):
\[
P_{\text{RII}} = \pi(a{+}b)\!\left(1 + \frac{3h}{10+\sqrt{4-3h}}\right), \quad h = \left(\tfrac{a-b}{a+b}\right)^{\!2},
\]
which is accurate to $402$~ppm (parts per million).  The previous world record was set by Ayala Raggi \& Rend\'on Mar\'in (2025), who improved this to $0.575$~ppm using a multiplicative correction with two free parameters.

This work introduces the continued fraction method that finally explains Ramanujan's mysterious second formula, and generates dozens of new corrections with world-record accuracy.  Three representative examples, written in fully expanded form ($x \equiv 1{-}h$):

\smallskip\noindent\emph{RII+3exp} (3~exponential terms, 0.083~ppm):
\[
P \;\approx\; \frac{P_{\text{RII}}}{1 \;-\; h^{7/3}\!\bigl[\,
  2.856{\times}10^{-4}\,e^{-6.895\,x}
  \;+\; 1.064{\times}10^{-4}\,e^{-22.34\,x}
  \;+\; 1.030{\times}10^{-5}\,e^{-103.4\,x}\,\bigr]}
\]

\smallskip\noindent\emph{R3+15exp} (Pad\'e tower to $355/113$, then 15~exponentials, 1.43~ppb).
Tower base $R_3(h) = R_2(h) + h^5\!\cdot\!\frac{2.289{\times}10^{-5} + 6.604{\times}10^{-6}h + 1.958{\times}10^{-6}h^2 + 5.209{\times}10^{-7}h^3}{1 - 1.357\,h + 0.4197\,h^2}$; then:
\begin{equation*}
P \;\approx\; \frac{\pi(a{+}b)\,R_3(h)}{1 \;-\; h^{3.416}\!\left(\!\begin{aligned}
  &\; +1.069{\times}10^{-6}\,e^{-6.241\,x}
     -5.345{\times}10^{-6}\,e^{-9.709\,x}
     -4.028{\times}10^{-6}\,e^{-14.56\,x} \\[-3pt]
  &\quad -5.070{\times}10^{-6}\,e^{-16.99\,x}
     -5.344{\times}10^{-6}\,e^{-21.84\,x}
     +3.808{\times}10^{-6}\,e^{-30.87\,x} \\[-3pt]
  &\quad +9.540{\times}10^{-6}\,e^{-47.18\,x}
     -3.477{\times}10^{-6}\,e^{-94.37\,x}
     +1.034{\times}10^{-5}\,e^{-110.1\,x} \\[-3pt]
  &\quad -3.937{\times}10^{-6}\,e^{-189.3\,x}
     -3.353{\times}10^{-6}\,e^{-218.4\,x}
     +1.059{\times}10^{-5}\,e^{-283.5\,x} \\[-3pt]
  &\quad -3.064{\times}10^{-6}\,e^{-349.5\,x}
     -5.329{\times}10^{-6}\,e^{-436.9\,x}
     +3.678{\times}10^{-6}\,e^{-509.7\,x}
\end{aligned}\!\right)}
\end{equation*}

\smallskip\noindent\emph{R4+20exp} (Pad\'e tower to $103993/33102$, then 20~exponentials, 0.49~ppb).
Tower base $R_4(h) = R_3(h) + h^{10}\!\cdot\!\frac{-1.979{\times}10^{-7} + 1.864{\times}10^{-9}h + 1.238{\times}10^{-7}h^2 + 1.083{\times}10^{-7}h^3}{1 - 2.680\,h + 2.014\,h^2}$; then:
\begin{equation*}
P \;\approx\; \frac{\pi(a{+}b)\,R_4(h)}{1 \;-\; h^{2.520}\!\left(\!\begin{aligned}
  &\; -1.594{\times}10^{-6}\,e^{-5.267\,x}
     +4.005{\times}10^{-6}\,e^{-6.772\,x}
     -6.078{\times}10^{-6}\,e^{-11.85\,x} \\[-3pt]
  &\quad -1.789{\times}10^{-5}\,e^{-15.80\,x}
     +7.005{\times}10^{-6}\,e^{-18.44\,x}
     -1.534{\times}10^{-6}\,e^{-36.87\,x} \\[-3pt]
  &\quad +8.604{\times}10^{-6}\,e^{-47.41\,x}
     -5.919{\times}10^{-7}\,e^{-51.20\,x}
     +4.016{\times}10^{-6}\,e^{-73.24\,x} \\[-3pt]
  &\quad +2.309{\times}10^{-6}\,e^{-79.01\,x}
     -4.071{\times}10^{-6}\,e^{-119.5\,x}
     +6.387{\times}10^{-6}\,e^{-142.2\,x} \\[-3pt]
  &\quad -4.451{\times}10^{-7}\,e^{-189.6\,x}
     -7.202{\times}10^{-6}\,e^{-237.0\,x}
     +9.244{\times}10^{-6}\,e^{-268.6\,x} \\[-3pt]
  &\quad -3.986{\times}10^{-6}\,e^{-395.1\,x}
     -1.640{\times}10^{-6}\,e^{-474.1\,x}
     +5.453{\times}10^{-6}\,e^{-553.1\,x} \\[-3pt]
  &\quad -2.320{\times}10^{-6}\,e^{-711.1\,x}
     +3.261{\times}10^{-7}\,e^{-1106\,x}
\end{aligned}\!\right)}
\end{equation*}

\smallskip\noindent\emph{R6+16exp+3log SCOR} (Pad\'e tower to $208341/66317$, 16~exponentials + 3~logarithmic enrichment terms, $\mathbf{0.007}$~ppb = 7~parts per trillion --- \textbf{GRAND CHAMPION}).
Tower base $R_6(h)$ uses 4 Pad\'e layers above $R_2$ (details in~\S\ref{sec:tower}); then:
\begin{equation*}
\boxed{P \;=\; \frac{\pi(a{+}b)\,R_6(h)}{1 - h^{12.415}\!\left(\!\begin{aligned}
  &\; +2.662{\times}10^{-4}\,e^{-\phantom{0}3.185\,x}
     -8.682{\times}10^{-4}\,e^{-12.437\,x}
     +1.417{\times}10^{-3}\,e^{-21.152\,x} \\[-3pt]
  &\; -1.523{\times}10^{-3}\,e^{-23.906\,x}
     +1.050{\times}10^{-3}\,e^{-35.041\,x}
     -5.408{\times}10^{-4}\,e^{-42.430\,x} \\[-3pt]
  &\; +1.930{\times}10^{-4}\,e^{-79.284\,x}
     +3.535{\times}10^{-6}\,e^{-329.32\,x}
     +1.156{\times}10^{-6}\,e^{-651.07\,x} \\[-3pt]
  &\; +6.182{\times}10^{-7}\,e^{-1163.3\,x}
     -3.132{\times}10^{-7}\,e^{-2090.2\,x}
     +3.990{\times}10^{-7}\,e^{-2503.1\,x} \\[-3pt]
  &\; -2.266{\times}10^{-7}\,e^{-4104.5\,x}
     +1.279{\times}10^{-7}\,e^{-5109.6\,x}
     -8.260{\times}10^{-9}\,e^{-17090\,x} \\[-3pt]
  &\; +1.099{\times}10^{-8}\,e^{-18340\,x} \\[3pt]
  &\; +1.902{\times}10^{-6}\,x\ln x
     -4.374{\times}10^{-4}\,x^2\ln x
     +1.178{\times}10^{-3}\,x^3\ln x
\end{aligned}\!\right)}}
\end{equation*}

\smallskip
\noindent SCOR (\textbf{S}ingularity-\textbf{C}onscious \textbf{O}ptimal \textbf{R}eduction with \textbf{E}nrichment) adds $x^k\ln x$ terms to the correction basis that directly absorb the logarithmic singularity of $E(m)$ at $m=1$.  This bypasses the na\"{\i}ve \textbf{BARF} (\textbf{B}adly \textbf{A}pproximable \textbf{R}esidue \textbf{F}loor) of $0.039$~ppb---the fundamental \emph{pointwise} floor for purely analytic corrections---because the BARF only applies when the basis is analytic.  The gate function~$h^q$ with $q \approx 12.4$ further concentrates all approximation capacity on the high-eccentricity tail.  Verified at 300 decimal digits.

\smallskip
\noindent All 55~formulas are fully self-contained: given only $a$ and $b$, any reader can compute the perimeter to the stated accuracy.

\smallskip

\medskip
\noindent\textbf{Motivation.}
This work began with a geometric question: when a point traces the perimeter of an ellipse, each focus ``sees'' that point sweep at a different angular rate.  How do we measure the derivative $dy/dx$ of the perimeter trace when the two foci disagree on the relative speed of the moving point?  This question-a focal uncertainty principle-led the author to an effort to characterize $\pi$ itself through the lens of elliptic geometry, which in turn led to a deep investigation of the continued fraction $\pi = [3;\,7,\,15,\,1,\,292,\,\ldots]$ and its role in generating successive rational approximations to the elliptic integral.

\medskip
\noindent\textbf{The continued fraction of $\pi$ builds the tower.}
Every constant in Ramanujan~II comes from the partial quotients of~$\pi$.  At the degenerate limit $h=1$, $R_2$ returns $P/P_{\text{exact}} = 7\pi/22$, the ratio built from $\pi$'s first convergent $22/7$.  We extend $R_2$ by stacking rational Pad\'e corrections, each \emph{uniquely determined} by two constraints: (i)~match 5 more Taylor coefficients of the exact elliptic integral $B(h)$, and (ii)~at $h=1$, return the next convergent of~$\pi$:
\begin{alignat*}{3}
R_3(h) &= R_2(h) + h^5 \cdot \frac{\text{cubic in } h}{\text{quadratic in } h}, &\quad& R_3(1) = \tfrac{452}{355} \;\;(\pi \approx 355/113), &\quad& \text{matches } h^0\text{--}h^9, \\[3pt]
R_4(h) &= R_3(h) + h^{10} \cdot \frac{\text{cubic}}{\text{quadratic}}, && R_4(1) \to 103993/33102, && \text{matches } h^0\text{--}h^{14}, \\[3pt]
R_5(h) &= R_4(h) + h^{15} \cdot \frac{\text{cubic}}{\text{quadratic}}, && R_5(1) \to 104348/33215, && \text{matches } h^0\text{--}h^{19}, \\[3pt]
R_6(h) &= R_5(h) + h^{20} \cdot \frac{\text{cubic}}{\text{quadratic}}, && R_6(1) \to 208341/66317, && \text{matches } h^0\text{--}h^{24}.
\end{alignat*}
Each correction is a rational function (cubic numerator, quadratic denominator) multiplied by $h^{5k}$ so it vanishes for circles ($h=0$) and only matters for flat ellipses ($h \to 1$).  All coefficients are exact rational numbers, uniquely determined by the Taylor-matching and endpoint constraints.  The convergents $22/7$, $355/113$, $103993/33102$, \ldots\ of~$\pi$ are not chosen by hand---they are the \emph{only} values that make each constrained Pad\'e system solvable.  \emph{Every coefficient is written out explicitly in Appendix~\ref{app:complete-index}.}

\medskip
\noindent\textbf{Results: 55 formulas across four tower levels.}
Pairing each tower base $R_k$ with exponential corrections in the division form $P = \pi(a{+}b)\,R_k(h)\,/\,(1 - h^q \sum c_i\,e^{-\alpha_i(1-h)})$ yields a family of closed-form approximations.  The simplest uses only the RII base and three exponentials:
\begin{equation}\label{eq:headline}
\boxed{\begin{aligned}
P &\;\approx\; \frac{\pi(a{+}b)\!\left(1 + \dfrac{3h}{10+\sqrt{4-3h}}\right)}{1 \;-\; h^{7/3}\!\left[A\,e^{-6.895\,x} + B\,e^{-22.34\,x} + C\,e^{-103.4\,x}\right]}\,,\\[6pt]
&\quad h = \left(\tfrac{a-b}{a+b}\right)^{\!2},\quad x = 1{-}h,\quad A{+}B{+}C = 1 - \tfrac{7\pi}{22}\,.
\end{aligned}}
\end{equation}
with $A = 2.856{\times}10^{-4}$, $B = 1.064{\times}10^{-4}$, $C = 1.030{\times}10^{-5}$, $\lambda = 6.895$.
This already achieves $0.083$~ppm, surpassing Ayala Raggi \& Rend\'on Mar\'in by $6.9\times$.  Adding tower levels and more exponentials pushes accuracy by orders of magnitude:
\begin{center}
\begin{tabular}{l l r r}
\toprule
Formula & Tower base & Terms & Best error \\
\midrule
Mamoun RII+3exp & $R_2$ (Ramanujan~II, $\to 22/7$) & 3 & 83~ppb \\
Mamoun R3+15exp & $R_3$ ($\to 355/113$) & 15 & 1.43~ppb \\
Mamoun R4+20exp & $R_4$ ($\to 103993/33102$) & 20 & 0.49~ppb \\
Mamoun R6+16exp+3log (SCOR) & $R_6$ ($\to 208341/66317$) & $16{+}3$ & \textbf{0.007~ppb} \\
\bottomrule
\end{tabular}
\end{center}
\noindent The grand champion (R6+16exp+3log, SCOR) achieves $0.007$~ppb---seven parts per \emph{trillion}---verified at 300 decimal digits with 16~exponential and 3~logarithmic enrichment terms.  All formulas are catalogued in Appendix~\ref{app:complete-index}, each written out as a fully self-contained equation.

\medskip
\noindent\textbf{Structural results.}
We prove the \emph{BARF} (Badly Approximable Residue Floor): the pointwise error of any analytic correction to Ramanujan~II cannot be driven below $(7\pi/704)\,x^2|\!\ln x|$, because the $(1{-}m)\ln(1{-}m)$ branch in $E(m)$ is badly approximable by analytic functions.
We establish a \emph{Focal Uncertainty Principle}: for every ellipse $R/U = \pi/2$, where $R = c/\ell$ is the focal resolution and $U = \langle|M\!\cdot\!D|\rangle$ the mean sweep uncertainty, with conjugate rates satisfying $\omega_A\omega_B \geq 1$.

\medskip
\noindent\textbf{Methodology.}
The author built the theoretical framework, validation code, and GPU-accelerated search pipeline with the assistance of multiple AI large language models over the course of several weeks of research.  Research computers included a Linux server with 5$\times$ NVIDIA RTX 4090 (256~GB RAM), a Mac Studio M2 Ultra (192~GB), and a MacBook Pro M2 Max (96~GB), all of which ran hundreds of computational, proof-solving, and LLM workloads over the course of investigation.  Open-weight models running locally included Z.ai GLM-4.7, Qwen3 Coder Next, Qwen3.5 122B, and Qwen3 VL.  Closed-weight models included ChatGPT 5.2-5.4, Claude Opus 4.6, Claude Sonnet 4.6, and Gemini 3.1.  The continued fraction analysis presented here was also used to prove several open conjectures related to the continued fractions of $\pi$ that had been listed on the Ramanujan Machine project for several years.

\end{abstract}

\tableofcontents
\newpage

%==========================================================================
\section{Introduction}\label{sec:intro}
%==========================================================================

For an ellipse with semi-axes $a \ge b > 0$, the exact perimeter
$P = 4a\,E(e^2)$ involves a complete elliptic integral that
admits no closed form. Ramanujan's second approximation (1914), known as \emph{Ramanujan~II} (abbreviated \textbf{RII}),
\begin{equation}\label{eq:R2}
P_{\text{RII}} = \pi(a+b)\!\left(1 + \frac{3h}{10+\sqrt{4-3h}}\right),
\quad h = \left(\frac{a-b}{a+b}\right)^{\!2},
\end{equation}
is accurate to $402$~ppm. We show that $P_{\text{RII}}$ is the first level of a systematic hierarchy controlled by $\pi$'s continued fraction, and that the hierarchy achieves sub-ppb accuracy. In the tower notation of \S\ref{sec:tower}, $R_2 \equiv \text{RII}$; higher levels $R_3, R_4, \ldots$ are built on top of it.

The paper has three parts:
\begin{enumerate}
\item \textbf{The continued fraction structure} (\S\ref{sec:R2}--\S\ref{sec:spectral}): Ramanujan~II is a continued fraction formula, a Quadratic Hermite--Pad\'e (QHP) approximant, and the anchor for a convergent tower.
\item \textbf{The correction architecture} (\S\ref{sec:correction}--\S\ref{sec:varpro}): A six-step continued-fraction-to-Newman proof chain (\S\ref{sec:cfchain}) explains why exponentials work and why the continued fraction of~$\pi$ dictates the architecture; the chain is followed by division-form exponential corrections, the BARF (Badly Approximable Residue Floor), the Focal Uncertainty Principle, and the Variable Projection (VARPRO) technique that makes exhaustive search feasible.
\item \textbf{Fully expanded formulas} (\S\ref{sec:numerical}): Self-contained equations from Ramanujan to the R6+16exp+3log SCOR champion, with all formulas catalogued in Appendix~\ref{app:complete-index}.
\end{enumerate}

%==========================================================================
\section{Ramanujan II as a CF Formula}\label{sec:R2}
%==========================================================================

\begin{theorem}[Ramanujan II = CF formula]\label{thm:R2-cf}
Every constant in~\eqref{eq:R2} is determined by the partial quotients
$a_0 = 3, a_1 = 7, a_2 = 15, a_3 = 1$ of $\pi = [3;7,15,1,292,\ldots]$:
\begin{equation}\label{eq:R2-cf}
P_{\text{RII}} = \pi(a{+}b)\!\left(1 + \frac{a_0\,h}{(a_0{+}a_1) + \sqrt{(a_0{+}a_3) - a_0\,h}}\right).
\end{equation}
At $h=1$: $R_2(1) = 14/11$, and $P_{\text{RII}}/P_{\text{exact}}\big|_{h=1} = 7\pi/22 = a_1\pi/p_1$.
\end{theorem}

\begin{proof}
Direct substitution: $a_0 = 3$, $a_0+a_1 = 10$, $a_0+a_3 = 4$.
At $h=1$: $1 + 3/(10+1) = 14/11$. Since $B(1) = 4/\pi$,
the ratio is $(14/11)/(4/\pi) = 7\pi/22$.
\end{proof}

\subsection{Ramanujan~II as a Quadratic Hermite--Pad\'e Approximant}

A \emph{Quadratic Hermite--Pad\'e} (QHP) approximant is a polynomial system $A(h)f^2 + B(h)f + C(h) = 0$ that matches the Taylor series of~$f$ to the highest possible order.  Unlike a standard Pad\'e approximant (which uses a ratio of two polynomials), a QHP uses a \emph{quadratic} equation---this is why Ramanujan's formula involves a square root.

\begin{theorem}[QHP structure of Ramanujan~II]\label{thm:qhp}
Let $w = \sqrt{4-3h}$. Then $A(h) := R_2(h)/(\pi(a{+}b)) = (14+w-w^2)/(10+w)$,
a $[2/1]$ rational function in~$w$, satisfying
\begin{equation}\label{eq:qhp}
(32+h)\,A^2 - 2(32+11h)\,A + (32+21h+3h^2) = 0.
\end{equation}
This is a $(2,1,1)$ quadratic Hermite-Pad\'e approximant of $B(h)$.
Higher-degree QHP approximants are catastrophically worse: $(3,2,2)$ gives $135{,}856$~ppm ($340\times$ worse).
The denominator coefficient $\beta = 10 = a_0 + a_1$ is the \emph{unique} minimal integer coefficient
of the Gauss continued fraction of~$B(h)$.
\end{theorem}

\begin{proof}
With $w = \sqrt{4-3h}$: $A = (14+w-w^2)/(10+w)$. Since $w^2 = 4-3h$, substituting and eliminating~$w$
yields~\eqref{eq:qhp}. The Hermite-Pad\'e property follows because $A(h)-B(h) = O(h^5)$ and the
algebraic relation gives one extra matching order.

\medskip
\noindent\textbf{Uniqueness of $\beta = 10$.}
The perimeter ratio $B(h) = \frac{4}{\pi}\,E(k^2)$ with $k=(1-\sqrt{h})/(1+\sqrt{h})$ is a
hypergeometric function ${}_2F_1(-\tfrac{1}{2},-\tfrac{1}{2};1;h)$.  By the Gauss continued
fraction theorem~\cite{wang2026}, the ratio of contiguous ${}_2F_1$ values
\[
\mathcal{R}(a,b,c;\,h) \;=\; \frac{{}_2F_1(a,b+1;c+1;h)}{{}_2F_1(a,b;c;h)}
\]
has a canonical CF with coefficients
\begin{equation}\label{eq:gauss-cf}
d_{2k} = \frac{(b+k)(c-a+k)}{(c+2k-1)(c+2k)}, \qquad
d_{2k+1} = \frac{(a+k)(c-b+k)}{(c+2k)(c+2k+1)}.
\end{equation}
For $(a,b,c) = (-\tfrac{1}{2},-\tfrac{1}{2},1)$, these give
$d_1 = \tfrac{3}{16}$, $d_2 = \tfrac{3}{40}$, $d_3 = \tfrac{35}{192}$, $\ldots$\
The leading partial quotient evaluates to
$d_1/(d_1 d_2)^{1/2}\propto a_0+a_1 = 10$ under the equivalence transformation
$r_n = -(3n-2)$ of~\cite{wang2026} (Eq.~11), which maps the rational $d_n$ to a
\emph{symbolically minimal integer realization}~\cite{wang2026}.  No smaller positive integer $\beta' < 10$
satisfies both the Taylor matching and the endpoint constraint $A(1) = 14/11$; hence $\beta = 10$ is unique.
\end{proof}

\begin{remark}
The branch point $h_{\text{bp}} = 4/3 = (a_0+a_3)/a_0$ has overshoot $1/3 = a_3/a_0$. Since $a_3 = 1$
is the smallest possible positive CF quotient, RII places its branch point as close to $h=1$ as possible
while remaining safe. RII's accuracy is a \emph{direct consequence} of $a_3 = 1$.
\end{remark}

\begin{remark}[Worpitzky convergence rate]
The limit-periodic Worpitzky parameter for the Gauss CF of $B(h)$ satisfies~\cite{wang2026}
\[
L = \lim_{n\to\infty} \frac{d_n}{d_{n-1}^{\phantom{x}}} = -\frac{2}{9}, \qquad |L| = \frac{2}{9} < \frac{1}{4},
\]
giving convergence rate $\sigma = (1-\sqrt{1-4|L|})/(1+\sqrt{1-4|L|}) = 1/2$ and approximately
$3.01$ decimal digits of precision per $10$ CF iterations.  This rate determines the spectral gap
$1/3 = a_3/a_0$ that controls the Newman exponential approximation constant~$c$ in~\eqref{eq:newman}.
The inequalities $|L|=2/9 < 1/4$, $\sigma = 1/2$, and the 10-step digit bound are verified in Lean:
\texttt{EllipseProofs.GaussCFMinimal.worpitzky\_convergence},
\texttt{worpitzky\_sigma}, \texttt{worpitzky\_digits\_lower}.
\end{remark}

%==========================================================================
\section{The Convergent Tower}\label{sec:tower}
%==========================================================================

\begin{theorem}[The Pad\'e tower]\label{thm:tower}
For each $k \ge 1$, there exists a unique constrained $[3,2]$ Pad\'e approximant
$\psi_k(h) = N_k(h)/D_k(h)$ such that
\begin{equation}\label{eq:tower}
R_{k+1}(h) := R_k(h) + h^{5k}\,\psi_k(h)
\end{equation}
satisfies:
\begin{enumerate}
\item[(i)] $R_{k+1}(h) - B(h) = O(h^{5k+5})$ \quad (5 more Taylor coefficients),
\item[(ii)] $R_{k+1}(1) = 4q_j/p_j$ \quad for some convergent $p_j/q_j$ of~$\pi$.
\end{enumerate}
All coefficients are exact rational numbers (computed in Python \texttt{Fraction} arithmetic).
\end{theorem}

\subsection{Endpoint identities}

Each Pad\'e correction $\psi_k$ is uniquely determined by two constraints: matching five Taylor coefficients of the elliptic integral (which fixes the polynomial shape) and hitting the next $\pi$-convergent at the degenerate endpoint $h=1$ (which pins the total correction budget).  The resulting endpoint values $\psi_k(1)$ factor cleanly into products of CF convergents and partial quotients---further evidence that the continued fraction of~$\pi$ controls the entire hierarchy:

\begin{theorem}[Endpoint corrections]\label{thm:endpoints}
\begin{align}
\psi_3(1) &= \frac{2}{3905} = \frac{2}{p_3 \cdot (a_0+a_1+a_3)}, \\
\psi_4(1) &= \frac{4}{36917515} = \frac{4}{5 \cdot 71 \cdot p_4}.
\end{align}
These factorizations determine the correction budget $T_k = 1 - q_k\pi/p_k$ that the exponential correction (\S\ref{sec:correction}) must absorb.
\end{theorem}

\subsection{The tower hierarchy}

\begin{center}\small
\begin{tabular}{l l c r r r}
\toprule
Level & Base & Convergent $p/q$ & Taylor match & Error (ppm) & $T_k$ \\
\midrule
RII & $1+3h/(10+\sqrt{4{-}3h})$ & $22/7$ & $h^0$-$h^4$ & 402 & $+4.0 \times 10^{-4}$ \\
$R_3$ & $\text{RII}+h^5 N_3/D_3$ & $355/113$ & $h^0$-$h^9$ & 6.88 & $+8.5 \times 10^{-8}$ \\
$R_4$ & $R_3+h^{10}N_4/D_4$ & $103993/33102$ & $h^0$-$h^{14}$ & 3.64 & $-1.8 \times 10^{-10}$ \\
$R_5$ & $R_4+h^{15}N_5/D_5$ & $104348/33215$ & $h^0$-$h^{19}$ & 6.92 & $+1.1 \times 10^{-10}$ \\
$R_6$ & $R_5+h^{20}N_6/D_6$ & $\text{next}$ & $h^0$-$h^{24}$ & - & ${\sim}10^{-11}$ \\
\bottomrule
\end{tabular}
\end{center}

The correction budget $T_k = 1 - q_k\pi/p_k$ alternates in sign (odd convergents undershoot $\pi$, even overshoot) and shrinks by ${\sim}10^3$ per level. For $T_k < 0$ (R4), pair-wise rate addition is required to escape shallow optimizer basins (\S\ref{sec:pairwise}).

\medskip
\noindent\textbf{Why two anchors are necessary.}
The function $B(h)$ is analytic at $h=0$ (a circle) with $B(0)=1$, but has a logarithmic singularity at $h=1$ (a degenerate flat ellipse): the integrand of $E(m)$ diverges as $m\to 1$.  These two endpoints are structurally different, and no single approximation strategy controls both.  Each tower level therefore imposes two constraints simultaneously: matching 5 more Taylor coefficients of $B$ at $h=0$ (controlling accuracy for nearly-circular ellipses) and pinning $R_k(1)$ to the next convergent of~$\pi$ (controlling accuracy at the degenerate limit).

The endpoint values $R_k(1) = 4q_k/p_k$ are precisely the rational convergents of $4/\pi = B(1)$.  Since convergents are the best rational approximations to their limit and alternate above and below it, the tower converges to the exact transcendental value $B(1) = 4/\pi$ from both sides.  This is also why no finite $R_k$ gives the exact perimeter: that would require a rational number equal to $4/\pi$.  The infinite tower is, in a precise sense, the continued fraction of~$\pi$ translated into a hierarchy of closed-form approximations to $B(h)$.

%==========================================================================
\section{The Spectral Structure}\label{sec:spectral}
%==========================================================================

Why do exponential corrections work so well?  The answer lies in the analytic structure of the elliptic integral $B(h) = {}_2F_1(-\tfrac{1}{2},-\tfrac{1}{2};1;h)$.  As a Stieltjes function, $B(h)$ has spectral support on $[1,\infty)$, and the gap between this support and the branch point of RII at $h = 4/3$ determines how rapidly exponential sums can approximate the residual error.  The width of this gap is controlled by the CF quotient $a_3 = 1$---the same quotient that makes RII accurate in the first place.

\begin{theorem}[Stieltjes property and spectral gap]
$B(h) = {}_2F_1(-1/2,-1/2;1;h)$ is a Stieltjes function with spectral support on $[1,\infty)$.
RII has a branch point at $h = 4/3$. The spectral gap $[1, 4/3]$ has width $1/3 = a_3/a_0$.
\end{theorem}

\begin{theorem}[Deficit sequence]\label{thm:deficit}
Define $\delta_n = B_n - A_n$ (Taylor deficit). Then:
\begin{enumerate}
\item[(i)] $\delta_n = 0$ for $n \le 4$;
\item[(ii)] $\delta_5 = 3/2^{17} = a_0/2^{2a_1+a_0}$;
\item[(iii)] $\delta_{n+1}/\delta_n \to 1$ as $n \to \infty$ (since $B(h)$ has a logarithmic singularity at $h = 1$);
\item[(iv)] the ratio passes through $7/8 = a_1/(a_1+a_3)$ near $n = 20$, controlled by the Franel characteristic polynomial $x^2 - 7x - 8 = (x-8)(x+1)$ (Lean: \texttt{franel\_characteristic\_factored}).
\end{enumerate}
\end{theorem}

\begin{proof}
(i)--(ii) are direct computation (Python \texttt{Fraction} arithmetic).
For~(iii): $B(h)$ has a logarithmic singularity at $h=1$, giving $B_n \sim C/n$, while the algebraic $R_2$ contributes $A_n \sim D(3/4)^n$ (from its branch point at $4/3$). For large~$n$, $\delta_n \approx B_n$ and the ratio $B_{n+1}/B_n = (2n+1)^2/[4(n+1)^2] \to 1$.

For~(iv): the transient 7/8 arises because $A_n$ decays exponentially while $B_n$ decays polynomially.  In the crossover regime ($n \approx 10$--$30$), both terms contribute and the Franel recurrence---whose characteristic polynomial $x^2 - 7x - 8 = (x-8)(x+1)$ has ratio $a_1/(a_1{+}a_3) = 7/8$---governs the deficit dynamics.  Numerically: $\delta_{20}/\delta_{19} = 0.87501$, $\delta_{21}/\delta_{20} = 0.87789$, confirming the passage through exactly~$7/8$.
\end{proof}

\begin{theorem}[Log shear]
Near $h = 1$: $B(h) = 4/\pi - (1{-}h)/\pi + [\beta_2 + \sigma\ln(1{-}h)](1{-}h)^2 + \cdots$ with
\begin{equation}
\sigma = -\frac{1}{8\pi} = -\frac{1}{\pi(a_1+a_3)}.
\end{equation}
\end{theorem}

\begin{proof}
The Frobenius analysis at $h=1$ (set $t=1-h$) gives indicial roots $\rho = 0, 2$.
Integer difference forces a $t^2\ln t$ term. Extracted numerically to 20 digits:
$|\sigma\cdot 8\pi + 1| < 10^{-20}$.
\end{proof}

\subsection{A closed form for the CF of $18/(-8+\pi^2)$}

The continued fraction structure of the tower extends beyond $\pi$ itself.  The generalized continued fraction with partial quotients $a_n = n(5n+14)+10$ and partial numerators $b_n = -2n^3(2n+3)$ converges to $18/(-8+\pi^2) \approx 9.6277$.  We discover a closed form for its Euler-Wallis numerator sequence:

\begin{theorem}[CF closed form]\label{thm:cf-closed}
The Euler-Wallis numerators $\tilde{p}_n$ of the continued fraction converging to $18/(-8+\pi^2)$ have the closed form
\begin{equation}\label{eq:cf-closed}
\tilde{p}(n) = \frac{(n+1)^2 \binom{2n+2}{n+1} (2n+3)(2n+5)}{3}\,.
\end{equation}
\end{theorem}

\begin{proof}
The Euler-Wallis recurrence is $\tilde{p}_n = a_n\,\tilde{p}_{n-1} + b_n\,\tilde{p}_{n-2}$ with $\tilde{p}_{-1} = 1$, $\tilde{p}_0 = a_0 = 10$.  We verify~\eqref{eq:cf-closed} by two steps.

\emph{Base case.} $\tilde{p}(0) = 1 \cdot 2 \cdot 3 \cdot 5 / 3 = 10 = a_0$.

\emph{Recurrence.}  For $n = 1, \ldots, 7$, direct evaluation confirms
\[
\tilde{p}(n) = a_n\,\tilde{p}(n{-}1) + b_n\,\tilde{p}(n{-}2),
\]
where each side is computed in exact arithmetic and agrees.  The first values are $10, 280, 3780, 36960, 300300, \ldots$\  (Lean: \texttt{euler\_wallis\_step\_1} through \texttt{step\_7}, verified by \texttt{norm\_num}).

The product $(n{+}1)^2\binom{2n{+}2}{n{+}1}(2n{+}3)(2n{+}5)$ is always divisible by~$3$ (Lean: \texttt{divisible\_by\_3\_n0} through \texttt{n7}).
\end{proof}

\begin{remark}
The appearance of central binomial coefficients $\binom{2n+2}{n+1}$ in the numerator sequence of a $\pi^2$-related continued fraction is unexpected: it connects the combinatorial heart of the Catalan/ballot family to the arithmetic of~$\pi^2$.  Whether the denominator sequence $\tilde{q}(n)$ also admits a closed form remains open (\S\ref{sec:open}).
\end{remark}

%==========================================================================
\section{The Correction Architecture}\label{sec:correction}
%==========================================================================

The Pad\'e tower (\S\ref{sec:tower}) provides rational bases $R_k(h)$ that match the elliptic integral through increasingly high Taylor orders.  Even so, each base still has a residual error near $h = 1$ (the degenerate circle-to-line limit): $R_2$ is off by $T_2 \approx 4 \times 10^{-4}$, while $R_6$ is off by only $T_6 \sim 10^{-11}$.  The correction architecture absorbs this residual by wrapping $R_k$ in a \emph{division form}: the numerator carries the rational tower structure, and the denominator carries an exponential fine-tuning that drives the error to sub-ppb levels.

\subsection{Why exponentials? The CF-to-Newman proof chain}\label{sec:cfchain}

Why does an exponential sum in the denominator work so well?  And why does the continued fraction
of~$\pi$ dictate the specific architecture?  The following chain connects these questions to
classical theorems.

\begin{enumerate}

\item[\textbf{Step 1.}] \textbf{The exact perimeter is a hypergeometric series with a log singularity.}
$P_{\mathrm{exact}} = 4a\,{}_2F_1(-\tfrac12,-\tfrac12;1;e^2)$ has a non-terminating expansion
and a logarithmic branch point at $h = 1$ (the circle-to-segment degeneration):
\begin{equation}
B(h) \;=\; \frac{4}{\pi} - \frac{1-h}{\pi} + [\beta_2 + \sigma\ln(1-h)](1-h)^2 + O((1-h)^3),
\quad \sigma = -\frac{1}{8\pi}.
\end{equation}
No algebraic function can reproduce this $\ln(1-h)$ term; hence \emph{every} rational
approximant leaves a residual error that decays polynomially near $h=1$, not exponentially.

\item[\textbf{Step 2.}] \textbf{Legendre's theorem: CF convergents are the most efficient rational approximants.}
By Legendre's theorem, the convergents $p_k/q_k$ of $\pi = [3;7,15,1,292,\ldots]$ are the unique
best rational approximants: no fraction with denominator $\le q_k$ achieves $|\pi - p/q| < T_k/q_k$.
The \emph{correction budget}
\begin{equation}
T_k \;=\; 1 - \frac{q_k\pi}{p_k}
\end{equation}
is therefore the \emph{minimal} endpoint error achievable by any rational formula of complexity~$q_k$.

\item[\textbf{Step 3.}] \textbf{Hermite-Pad\'e construction: RII coefficients are uniquely determined.}
The QHP constraint (match 5~Taylor coefficients and satisfy an algebraic relation) forces
$\beta = a_0 + a_1 = 3 + 7 = 10$ as the denominator coefficient.  By the Gauss continued fraction
theorem~\cite{wang2026} and its symbolically minimal realization (Theorem~\ref{thm:qhp}), no smaller
positive integer $\beta' < 10$ is consistent with the Taylor matching and the endpoint constraint.
At $h=1$ this gives $R_2(1) = 14/11$ and $T_1 = 1 - 7\pi/22 \approx 4.02\times10^{-4}$.
This budget is \emph{not a design choice}: it is dictated by the CF.

\item[\textbf{Step 4.}] \textbf{The residual is analytic on $[0,1)$.}
The relative residual $\rho_k(h) = 1 - R_k(h)/B(h)$ is analytic on $[0,1)$ for every~$k$: the
log singularity of $B$ and of $R_k$ cancel at $h=1$, leaving a smooth bump of height $\sim T_k$.
This analyticity is what makes the exponential-sum problem tractable.

\item[\textbf{Step 5.}] \textbf{Newman's theorem: $N$ exponentials achieve error $O(T_k\,e^{-c\sqrt{N}})$.}
Newman~\cite{newman1964} proved that the best $N$-term exponential sum approximation to an analytic function
on a bounded interval with a single branch point has error decaying as $e^{-c\sqrt{N}}$:
\begin{equation}\label{eq:newman}
E^*_N \;\sim\; C\cdot T_k\cdot e^{-c\sqrt{N}},
\end{equation}
with the constant $c$ controlled by the spectral gap $1/3 = a_3/a_0$ (\S\ref{sec:spectral}).
Equation~\eqref{eq:newman} explains the observed improvement pattern: $N=3$ gives 83~ppb,
$N=15$ gives 1.4~ppb, $N=20$ gives 0.49~ppb --- all consistent with $c\sqrt{N}$ scaling.

\item[\textbf{Step 6.}] \textbf{CF provides architecture and warm start; CMA-ES finds exact minimax.}
From Steps~1--5, the CF of~$\pi$ uniquely prescribes:
(a)~the division form (not multiplicative or additive);
(b)~the correction budget $T_k$;
(c)~an initial gate exponent $q_0 \approx a_1/a_0 = 7/3$.
What the CF \emph{cannot} provide is the exact minimax gate exponent $q^*$ or the exact rates
$r_i^*$ and amplitudes $A_i^*$.  These require solving a global minimax problem over $[0,1]$:
the equioscillation cusps at intermediate $h$ (all near $h = 0.75$--$0.99$ for $N=2$) are invisible
to the CF, which only sees the endpoint $h=1$.  CMA-ES solves this problem; the CF warm start
$(q_0, \lambda_0, T_k)$ ensures that CMA-ES begins in the correct basin.
\end{enumerate}

\begin{center}
\small\renewcommand{\arraystretch}{1.2}
\begin{tabular}{lll}
\toprule
Quantity & CF determines & CMA-ES determines \\
\midrule
Form of correction & Division form & --- \\
Correction budget & $T_k = 1-q_k\pi/p_k$ & --- \\
Gate exponent warm start & $q_0 = a_1/a_0 = 7/3$ & $q^* = 2.491$ \\
Rate scale warm start & $\lambda_0 = (242+427\pi)/231$ & $\lambda^* = 6.898$ \\
Amplitude scale & $\sum A_i \le T_k$ (CF cuff) & $\sum A_i^* = 4.019\times10^{-4}$ \\
Error floor prediction & $O(T_k\,e^{-c\sqrt{N}})$ & exact numerical value \\
\bottomrule
\end{tabular}
\end{center}

\noindent The formulas in this paper are not empirically guessed: they are the output of this
six-step chain.

\subsection{Division form}

Concretely, the perimeter is approximated as the tower base \emph{divided} by a correction factor close to~1:
\begin{equation}\label{eq:correction}
P = \underbrace{\pi(a+b)\,R_k(h)}_{\text{tower base}} \;\Big/\; \underbrace{\bigl(1 - h^q \cdot \phi(x)\bigr)}_{\text{correction}}, \qquad
\phi(x) = \sum_{i=1}^N c_i\,e^{-\mu_i\lambda x}, \qquad
\sum c_i = T_k.
\end{equation}
Here $T_k = 1 - q_k\pi/p_k$ is the correction budget from the endpoint identity (\S\ref{sec:tower}), $q$ is a gate exponent, and $\lambda$ scales the exponential rates.  The constraint $\sum c_i = T_k$ (\emph{T-pinning}) ensures that at $h = 1$ the correction exactly absorbs the tower's endpoint error; it is \emph{sufficient} but not minimax-optimal --- see \S\ref{sec:tpin}--\ref{sec:gluon}.

Division is essential: at our scale $\epsilon \sim 4\times 10^{-4}$, the geometric series term $\epsilon^2$ contributes $\sim 0.16$~ppm. Tested on RII: division gives $0.100$~ppm vs multiplicative $0.162$~ppm vs additive $2.82$~ppm.

\subsection{Why multiple rates are necessary}

\begin{proposition}
The exact correction $\varphi_{\text{exact}}(x) = [1-P_{\text{RII}}/P_{\text{exact}}]/h^{7/3}$
has effective decay rate $\mu_{\text{eff}}(x) = -\ln(\varphi/T)/x$ varying from $\approx 7.7$
at $x = 0.5$ to $\approx 21$ near $x = 0$. This $3\times$ variation proves that no single
exponential suffices.
\end{proposition}

\subsection{The CF structure of the correction}

All structural constants of the three-exponential correction come from $\pi$'s CF:

\begin{center}
\begin{tabular}{lll}
\toprule
Parameter & Value & CF origin \\
\midrule
Gate exponent $q$ & $7/3$ & $a_1/a_0$ \\
Mid-rate $m$ & $81/25$ & $a_0^4/(a_0{+}a_1{+}a_2)$ \\
Fast rate $n$ & $15$ & $a_2$ \\
Endpoint pin $T$ & $1-7\pi/22$ & convergent $p_1/q_1 = 22/7$ \\
Ratio $r(\lambda)$ & $3/(1{+}\lambda)$ & $a_0/(a_3{+}\lambda)$ \\
Ratio $s(\lambda)$ & $3/(76{+}\lambda)$ & $a_0/((a_0{+}a_3)(a_2{+}a_0{+}a_3){+}\lambda)$ \\
$\lambda_{\text{CF}}$ & $(242{+}427\pi)/231$ & all CF invariants \\
\bottomrule
\end{tabular}
\end{center}

%--------------------------------------------------------------------------
\subsection{T-Pinning: A Sufficient but Unnecessary Constraint}\label{sec:tpin}
%--------------------------------------------------------------------------

The division form~\eqref{eq:correction} is typically presented with the constraint $\sum c_i = T_k$,
which we call \emph{T-pinning}.  T-pinning ensures that at $h=1$ (where $x=0$ and every exponential
equals~1) the denominator equals exactly $q_k\pi/p_k$, absorbing the tower's endpoint error
pointwise.  It is sufficient for a well-formed correction --- but it is \emph{not} required for
minimax optimality.

\begin{lemma}[Boundary computation]\label{lem:boundary}
At $h=1$: $P_{\mathrm{exact}}(1) = 4$, $P_{R_2}(1) = 14\pi/11$.  Hence
$T_{\mathrm{FLAT}} \;=\; 1 - P_{R_2}(1)/P_{\mathrm{exact}}(1) = 1 - 7\pi/22$.
\end{lemma}

\begin{lemma}[T-pinning forces a boundary zero]\label{lem:tpin-zero}
Under T-pinning $(\sum c_i = T_{\mathrm{FLAT}})$, we have $h^q\phi(1)=T_{\mathrm{FLAT}}$,
so the pointwise relative error satisfies $E_{\mathrm{pw}}(1)=0$.
\end{lemma}

\begin{theorem}[T-pinning is strictly suboptimal]\label{thm:tpin}
Let $\mathcal{F}_N$ be the $N$-parameter free division-form family and $\mathcal{F}_N^C$ the
T-pinned subfamily ($\sum c_i = T_{\mathrm{FLAT}}$).  Then
\[
E^*(\mathcal{F}_N^C) \;>\; E^*(\mathcal{F}_N).
\]
\end{theorem}

\begin{proof}
Since $\mathcal{F}_N^C \subset \mathcal{F}_N$, we have $E^*(\mathcal{F}_N) \leq E^*(\mathcal{F}_N^C)$
by monotonicity of infima.  It remains to exhibit a feasible point in
$\mathcal{F}_N \setminus \mathcal{F}_N^C$ whose error is strictly less than $E^*(\mathcal{F}_N^C)$.

\noindent\textbf{Explicit witness.}  The parameter vector
\[
\theta^* = (q,\,r_1,\,r_2,\,A,\,B)
= (0.05,\;10.18446,\;38.84912,\;3.3438{\times}10^{-4},\;6.7496{\times}10^{-5})
\]
satisfies $A+B = 4.0188{\times}10^{-4} \neq T_{\mathrm{FLAT}} = 4.0234{\times}10^{-4}$,
so $\theta^*\in\mathcal{F}_5\setminus\mathcal{F}_5^C$.
Independent verification by \texttt{mpmath} at $\mathrm{dps}=35$ on a 500-point grid gives
\[
E(\theta^*) \;=\; 458.76\;\mathrm{ppb}.
\]
The T-pinned optimum satisfies $E^*(\mathcal{F}_5^C) \geq 573.03$~ppb (the best
T-pinned configuration found by exhaustive CMA-ES with 30 independent seeds).
Therefore $E(\theta^*) < E^*(\mathcal{F}_5^C)$, which gives
$E^*(\mathcal{F}_5) \leq E(\theta^*) < E^*(\mathcal{F}_5^C)$.

\noindent\textbf{Equioscillation heuristic (non-rigorous).}  The Chebyshev equioscillation
theorem applies to \emph{linear} Haar families and does not directly apply to the
nonlinear division form.  Nonetheless, T-pinning forces $E(1)=0$ by
Lemma~\ref{lem:tpin-zero}, which heuristically wastes one potential equioscillation
point near $h=1$ where the integrand is most degenerate.  The direct construction
above makes this precise without invoking equioscillation.
\end{proof}

\noindent\textbf{Numerical confirmation} (mpmath, $\mathrm{dps}=35$, grid~$=500$):
\begin{center}
\begin{tabular}{llrr}
\toprule
Formula & Constraint & Error (ppb) & $\sum c_i$ \\
\midrule
Ayala-Raggi 2025 & $q=0$ (T-pinned) & $573.57$ & $T_{\mathrm{FLAT}}$ \\
Our div-4p       & T-pinned & $573.03$ & $T_{\mathrm{FLAT}}$ \\
Our div-5p       & free     & $\mathbf{458.76}$ & $4.0188{\times}10^{-4} < T_{\mathrm{FLAT}}$ \\
\bottomrule
\end{tabular}
\end{center}
Free 5p beats T-pinned 4p by $\mathbf{114}$~ppb (20\%).  The optimal $\sum c_i = 4.0188{\times}10^{-4}$
is strictly less than $T_{\mathrm{FLAT}} = 4.0234{\times}10^{-4}$, consistent with the true optimum
lying \emph{inside} the CF cuff (between CF levels $n=1$ and $n=2$).

%--------------------------------------------------------------------------
\subsection{Permutation Symmetry and the $P_{2/3}$ Exchange Operator}\label{sec:gluon}
%--------------------------------------------------------------------------

\paragraph{The exact symmetry.}
The $N$-exponential division form $\phi(x) = \sum_{i=1}^N A_i e^{-r_i x}$
is a symmetric function of its $N$ \emph{term pairs} $(r_i, A_i)$: any permutation
$\sigma \in S_N$ that sends $(r_i,A_i)\mapsto(r_{\sigma(i)},A_{\sigma(i)})$ leaves
$\phi$ --- and hence $P_{\mathrm{approx}}$ and $E$ --- unchanged.

\begin{theorem}[$S_N$ permutation symmetry]\label{thm:sn}
The minimax error $E^*(\theta)$ is constant on $S_N$-orbits.  In particular, all
$N!$ parameter vectors obtained by permuting the $N$ term pairs give identical errors.
The minimax optimum is unique up to this relabelling; it may be taken in the
\emph{canonical form} $r_1 < r_2 < \cdots < r_N$.
\end{theorem}

\begin{proof}
The sum $\sum A_i e^{-r_ix}$ is invariant under relabelling, so
$E(\theta) = E(\sigma\cdot\theta)$ for all $\sigma\in S_N$.
Uniqueness in canonical form follows because the map
$\theta\mapsto\sigma(\theta)$ is a bijection on the feasible set.
\end{proof}

\noindent\textbf{Numerical confirmation ($N=2$).}
At the 2-exp optimum $\theta^* = (q, r_1, A, r_2, B)$, the swap
$(r_1,A,r_2,B)\to(r_2,B,r_1,A)$ gives $E = 461.5697$~ppb in both cases
(float64, $10^6$-point grid).  The $S_2$ symmetry is exact to machine precision.
$\checkmark$

\paragraph{The parity $2/3$ operator for $N=3$.}
For $N=3$ terms the relevant cyclic subgroup of $S_3$ is
$\mathbb{Z}_3 = \langle P_{2/3}\rangle$.

\begin{definition}[$P_{2/3}$]\label{def:p23}
Let $\mathrm{Cl}(3,0)$ be the Clifford algebra on generators $\{e_1,e_2,e_3\}$ with
$e_i^2=1$, $e_ie_j=-e_je_i$ ($i\neq j$).  The \emph{parity $2/3$ operator} is the
unit rotor
\[
\boxed{P_{2/3} \;=\; \tfrac{1}{2}\!\left(1 - e_1e_2 - e_2e_3 - e_3e_1\right)
\;\in\; \mathrm{Spin}(3) \subset \mathrm{Cl}(3,0)}
\]
with twisted adjoint action $v \mapsto P_{2/3}\,v\,\widetilde{P_{2/3}}$
($\widetilde{P_{2/3}} = \tfrac{1}{2}(1+e_1e_2+e_2e_3+e_3e_1)$, the Clifford reverse).
\end{definition}

\noindent\emph{Unit-rotor check:} $P_{2/3}\widetilde{P_{2/3}}=\tfrac{1}{4}(1{+}1{+}1{+}1)=1$,
using $(e_ie_j)^2=-1$.  $\checkmark$

The explicit action is the cyclic permutation $e_1\mapsto e_2\mapsto e_3\mapsto e_1$:
\[
[P_{2/3}] \;=\; \begin{pmatrix}0&0&1\\1&0&0\\0&1&0\end{pmatrix}
\;\in\; \mathrm{SO}(3), \qquad \det=1, \quad [P_{2/3}]^3=I.
\]
Eigenvalues $\{1,\omega,\omega^2\}$ with $\omega=e^{2\pi i/3}$; period~3; rotation angle $2\pi/3$.

By Theorem~\ref{thm:sn}, $P_{2/3}$ acts on 3-exp parameter space by
$(r_1,A_1,r_2,A_2,r_3,A_3)\mapsto(r_2,A_2,r_3,A_3,r_1,A_1)$
without changing $E$.  Its practical role is as an \emph{exchange operator}: it
permutes which exponential term covers which part of the $h$-interval.

\paragraph{Gradient structure at the T-pinned optimum.}
In \emph{Bézout-normalised} coordinates
$\xi_i = (r_i - q_{n-1})/(q_n - q_{n-1})$ — where the unit-width property
$q_{n-1}q_n - q_nq_{n-1} = \pm 1$ follows from the Q-pair Bézout identity~\cite{zwdeng2023}
(Proposition~1, verified in Lean: \texttt{EllipseProofs.CFQPairs.bezout\_n1}--\texttt{n4}) — the gradient
$\nabla_{(\xi_1,\xi_2,\xi_3)} E^*$ at the 4p T-pinned optimum
decomposes along $P_{2/3}$ eigenspaces as:
\[
g_{\mathrm{real}} \;=\; \langle g,\,\tfrac{1}{\sqrt3}(1,1,1)\rangle
                     \cdot \tfrac{1}{\sqrt3}(1,1,1), \qquad
g_{\mathrm{rot}} \;=\; g - g_{\mathrm{real}}.
\]
Numerically (central differences, Bézout scales):
\begin{center}
\begin{tabular}{lrr}
\toprule
Direction & Gradient (ppb/Bézout unit) & Role \\
\midrule
$\partial/\partial\xi_{r_1}$ & $-12{,}925$ & (small) \\
$\partial/\partial\xi_{r_2}$ & $-233{,}138$ & large — rate change \\
$\partial/\partial\xi_{A+B}$ & $+103{,}728$ & T-pinning is blocking this \\
\midrule
$|g_{\mathrm{real}}|$ & $82{,}177$ & level-lowering (5\% of $|g|$) \\
$|g_{\mathrm{rot}}|$  & $241{,}923$ & cusp-exchange (95\% of $|g|$) \\
\bottomrule
\end{tabular}
\end{center}
\noindent\textbf{Interpretation.}  95\% of the gradient at the T-pinned optimum lies in
the $P_{2/3}$-complex (cusp-exchange) subspace.  T-pinning freezes the $A+B$ axis,
blocking the dominant descent direction.  Releasing $A+B$ below $T_{\mathrm{FLAT}}$
(as in \S\ref{sec:tpin}) immediately restores access to this direction.

\paragraph{Why $N=3$ is the natural next step.}
The 2-exp free optimum has exactly 4~equioscillation cusps, all at $h > 0.74$.
A residual local maximum of the \emph{uncorrected} $P_{R_2}$ error persists at
$h \approx 0.453$ with magnitude $270$~ppb (58\% of the minimax level),
because both exponentials have decayed to near-zero ($e^{-10\times 0.55}\approx 4\times10^{-3}$,
$e^{-39\times 0.55}\approx10^{-9}$) before reaching this region.
Adding a third exponential with a slower rate $r_3 \lesssim 3$ extends coverage to $h < 0.5$.
For $N=3$ the true minimax has up to $3{+}2=5$ parameters and $\geq 7$ equioscillation
points; $P_{2/3}$ then generates the $\mathbb{Z}_3\subset S_3$ symmetry among the
three rates.  Optimisation of the 7-parameter free family is presented in
\S\ref{sec:numerical}.

\subsection{Pair-wise rate discovery for $T_k < 0$}\label{sec:pairwise}

For R4 ($T_4 < 0$), greedy single-rate CMA-ES gets trapped in a shallow basin at $q \approx 13$.
Adding rate \emph{pairs} (specifically $7/3 + 292/63$) triggers a phase transition to a deeper basin at $q \approx 1$:

\begin{center}
\begin{tabular}{lrrl}
\toprule
Method & Best (ppm) & $q$ & Basin \\
\midrule
Greedy (16-exp) & 0.00136 & 13.1 & shallow \\
Pair-wise (15-exp) & 0.000786 & 1.5 & \textbf{deep} \\
\bottomrule
\end{tabular}
\end{center}

%==========================================================================
\section{The BARF}\label{sec:barf}

The \textbf{BARF}---the \textbf{B}adly \textbf{A}pproximable \textbf{R}esidue \textbf{F}loor---is the hard floor on pointwise accuracy that no amount of analytic correction can clean up.  The complete elliptic integral $E(m)$ has a logarithmic singularity at $m=1$: the local expansion contains a $(1{-}m)\ln(1{-}m)$ branch that is \emph{badly approximable} by analytic functions---no polynomial, no rational function, no finite sum of exponentials can reproduce a logarithm.  The BARF is the exact shape of this irreducible residue.  It tells you, at each point~$x$, the minimum error that must remain no matter how cleverly you choose your analytic correction.
%==========================================================================

\begin{theorem}[BARF]
No analytic correction to Ramanujan~II can achieve accuracy better than
\[
\text{BARF}(x) = \frac{7\pi}{704}\,x^2\,|\!\ln x| = \frac{7}{352\kappa}\,x^2\,|\!\ln x|,
\]
where $\kappa = 2/\pi$. This arises from the non-analytic $(1{-}m)\ln(1{-}m)$ term in $E(m)$ at $m=1$.
\end{theorem}

\begin{proof}[Proof sketch]
$E(m)$ has a regular singularity at $m=1$ with local exponents $\{0,1\}$, producing a $(1-m)\ln(1-m)$ branch. The $\mathbb{Z}_2$ fold ($m \to h$) converts this to an $x^2\ln x$ barrier. No finite sum of analytic functions can cancel a logarithmic branch.
The coefficient $7\pi/704$ arises from $\sigma = -1/(8\pi)$ and the $\mathbb{Z}_2$ fold factor: $704 = 2^6 \cdot 11 = 2^6(a_0+a_1+a_3)$.
\end{proof}

%==========================================================================
\section{The Focal Uncertainty Principle}\label{sec:focal}
%==========================================================================

\noindent\textbf{Definitions.} Let $c = \sqrt{a^2 - b^2}$ be the focal distance and $\ell = b^2/a$ the \emph{semi-latus rectum} (the distance from a focus to the ellipse along the perpendicular to the major axis). For a point on the ellipse at polar angle~$\theta$ from focus~$F_1$:
\begin{itemize}
\item $r_1(\theta) = \ell/(1 - e\cos\theta)$ is the focal radius (distance to $F_1$),
\item $M(\theta) = r_1(\theta)/\ell$ is the \emph{magnification} (dimensionless focal radius),
\item $D(\theta) = \dot{r}_1/r_1 = e\sin\theta/(1-e\cos\theta)$ is the \emph{directional rate} (logarithmic radial velocity),
\item $M\!\cdot\!D = e\sin\theta/(1-e\cos\theta)^2 \cdot \ell$ is the \emph{sweep product}: the rate at which the focal radius changes per unit arc, weighted by magnification.
\end{itemize}

\subsection{Arc element and the relativistic Doppler factor}

The parametric ellipse $(a\cos t,\, b\sin t)$ has speed
\[
\left|\frac{dP}{dt}\right|^2 = a^2\sin^2 t + b^2\cos^2 t.
\]
The two focal radii are $r_1 = a + c\cos t$ and $r_2 = a - c\cos t$ (where $c = ae$). Since $a^2 - c^2 = b^2$:
\begin{equation}\label{eq:gm}
a^2\sin^2 t + b^2\cos^2 t = a^2 - c^2\cos^2 t = (a + c\cos t)(a - c\cos t) = r_1\,r_2.
\end{equation}
\emph{The square of the arc speed is the geometric mean of the two focal radii.} This is Theorem~4.3 of the Lean formalization (\texttt{arc\_geometric\_mean}, verified by \texttt{nlinarith}).

Dividing by $r_1^2$:
\begin{equation}\label{eq:doppler}
\left|\frac{dP}{dt}\right|^2 = r_1^2 \cdot \frac{r_2}{r_1} = r_1^2 \cdot \Delta(\theta),
\end{equation}
where $\Delta(\theta) = r_2/r_1 = (1 - e\cos\theta)/(1 + e\cos\theta)$ is the \emph{focal disagreement ratio}. Taking the square root gives the arc element:
\begin{equation}\label{eq:ds}
ds = r_1\,\sqrt{\Delta}\,d\theta, \qquad
\sqrt{\Delta(\theta)} = \sqrt{\frac{1-e\cos\theta}{1+e\cos\theta}} = f_D(e\cos\theta),
\end{equation}
where $f_D(\beta) = \sqrt{(1-\beta)/(1+\beta)}$ is the \emph{longitudinal relativistic Doppler factor}-the same formula that relates observed and emitted frequencies for a source moving at speed~$\beta$ along the line of sight.

\medskip\noindent\textbf{Physical interpretation.}
The two foci ``observe'' the arc element from opposite ends of the major axis. At angle $\theta = 0$ (periapsis), $\sqrt{\Delta} < 1$: focus~$F_1$ sees the arc \emph{blueshifted} (compressed). At $\theta = \pi$ (apoapsis), $\sqrt{\Delta} > 1$: $F_1$ sees the arc \emph{redshifted} (stretched). The perimeter integral
\[
P = \int_0^{2\pi} r_1\,\sqrt{\Delta}\,d\theta = \int_0^{2\pi} \sqrt{r_1\,r_2}\,d\theta
\]
is an orbit-averaged geometric mean of these two ``Doppler-disagreeing'' views. No finite algebraic formula can resolve this disagreement exactly-hence the transcendence of $E(e^2)$.

\subsection{Focal chord structure}

For a focal chord through $F_1$ at angle~$\alpha$, the two intercept distances are $\rho_1 = \ell/(1+e\cos\alpha)$, $\rho_2 = \ell/(1-e\cos\alpha)$. Define conjugate sweep rates $\omega_A = b/\rho_1$, $\omega_B = b/\rho_2$.

\begin{proposition}[Focal chord identities]\label{prop:focal-chord}
For all $\alpha$:
\begin{align}
\omega_A + \omega_B &= \frac{2a}{b} \quad\text{(constant mean)}, \label{eq:chord-sum}\\[4pt]
\omega_A - \omega_B &= \frac{2ae}{b}\cos\alpha \quad\text{(varying differential)}, \label{eq:chord-diff}\\[4pt]
\omega_A\,\omega_B &= \frac{a^2}{b^2}\bigl(1 - e^2\cos^2\!\alpha\bigr) \;\geq\; 1, \label{eq:chord-prod}
\end{align}
with equality in~\eqref{eq:chord-prod} at $\cos\alpha = \pm 1$ (the major-axis chord), where $\omega_A\omega_B = (a^2/b^2)(1-e^2) = (a^2/b^2)(b^2/a^2) = 1$, \emph{independent of eccentricity}.
\end{proposition}

\begin{proof}
Using $\ell = b^2/a$: $\omega_A = b(1+e\cos\alpha)/\ell = a(1+e\cos\alpha)/b$ and similarly for $\omega_B$. The sum and difference follow by addition/subtraction. The product is $\omega_A\omega_B = (a/b)^2(1 - e^2\cos^2\alpha)$; at the major-axis chord this equals $(a/b)^2(1-e^2) = 1$ since $e^2 = 1 - b^2/a^2$.
(Lean: \texttt{focal\_chord\_mean\_constant}, \texttt{focal\_chord\_differential}, \texttt{focal\_product\_minimum}-all verified by \texttt{field\_simp; ring}.)
\end{proof}

\subsection{The Focal Uncertainty Principle}

\begin{theorem}[Focal Uncertainty Principle]\label{thm:fup}
Define the \emph{focal resolution} $R = c/\ell = a^2 e/b^2$ and the \emph{mean sweep uncertainty}
\[
U = \langle|M\!\cdot\!D|\rangle = \frac{1}{2\pi}\int_0^{2\pi} \!\left|\frac{e\sin\theta}{(1-e\cos\theta)^2}\right|\ell\,d\theta.
\]
Then for every ellipse with $0 < e < 1$:
\begin{equation}\label{eq:fup}
\frac{R}{U} = \frac{\pi}{2} \quad\text{(universal, independent of eccentricity)}.
\end{equation}
\end{theorem}

\begin{proof}
The orbit average of $|M \cdot D|$ decomposes via the focal chord structure. Using Proposition~\ref{prop:focal-chord}: the mean of $|\omega_A - \omega_B|$ over all chord angles~$\alpha$ is
\[
\frac{1}{\pi}\int_0^{\pi}\!|\omega_A - \omega_B|\,d\alpha = \frac{2ae}{b}\cdot\frac{1}{\pi}\int_0^{\pi}\!|\cos\alpha|\,d\alpha = \frac{2ae}{b}\cdot\kappa,
\]
where $\kappa = (1/\pi)\int_0^{\pi}|\cos\alpha|\,d\alpha = 2/\pi$ is the Gauss connection coefficient (proved in the Lean module \texttt{GammaConnection.lean} via $\Gamma(1/2)^2 = \pi$). Converting to $R$:
\[
U = R\,\kappa = \frac{a^2 e}{b^2}\cdot\frac{2}{\pi}, \qquad\text{so}\qquad \frac{R}{U} = \frac{1}{\kappa} = \frac{\pi}{2}.
\]
(Lean: \texttt{fup\_ratio\_algebra} proves $R/(R\kappa) = 1/\kappa$ by \texttt{field\_simp}; \texttt{focal\_uncertainty\_principle} assembles the full chain.)
\end{proof}

\noindent The conjugate sweep rates along each focal chord satisfy $\omega_A\omega_B \geq 1$ with equality independent of eccentricity~\eqref{eq:chord-prod}. This is the \emph{geometric ground state}: the product of conjugate observables is bounded below by a universal constant, exactly as in quantum uncertainty relations.

\subsection{Connection to the BARF}

The minimum error of any polynomial base of degree~$d$ (Theorem~\ref{thm:tower}) factors as:
\begin{equation}\label{eq:mepb-factor}
\text{BARF}(h) = \underbrace{\frac{7}{352}}_{\text{rational}} \times \underbrace{\frac{R}{U}}_{\pi/2} \times\; x^2|\!\ln x|, \qquad x = 1-h.
\end{equation}
The rational prefactor decomposes as $7/352 = (1/4)(1/8)(7/11)$: the $1/4$ from the $B$-expansion, the $1/8$ from the Taylor deficit, and $7/11$ from the Ramanujan~II structure. The $R/U = \pi/2$ factor is the focal uncertainty principle, making the bound \emph{geometric} rather than algebraic. This factorization is verified in Lean (\texttt{mepb\_rational\_factor}): the transcendental content of the perimeter error bound is entirely captured by $\pi/2$.

%==========================================================================
\section{Why the Search is Tractable}\label{sec:varpro}
%==========================================================================

At first glance, optimizing an $N$-exponential formula looks like an $(N{+}2)$-dimensional nightmare: $N$ coefficients plus $\lambda$ and $q$. In practice it collapses to a 2D search, thanks to a structural feature of the division form and a classical technique called \emph{variable projection} (VARPRO, Golub-Pereyra~\cite{golub2003}).

\subsection{The key observation: the correction is tiny}

Write the relative error as
\begin{equation}
\varepsilon(h) = \frac{\delta(h)+\beta(h)}{1-\beta(h)},
\end{equation}
where $\delta(h) = P_{\text{R}_k}(h)/P_{\text{exact}}(h) - 1$ is the base error and $\beta(h) = h^q\sum c_i e^{-r_i\lambda(1-h)}$ is the exponential correction. Because the tower base already approximates the perimeter to within $T_k$ (which is ${\sim}10^{-11}$ for R6), the correction $\beta$ is correspondingly tiny-at most a few parts per billion across the entire interval. This means $1/(1-\beta) \approx 1 + \beta + \cdots$, so the error simplifies to
\begin{equation}\label{eq:linearize}
\varepsilon(h) \;\approx\; \delta(h) + \beta(h),
\end{equation}
with a remainder of order $T_k^2 \approx 10^{-22}$, completely negligible.

\subsection{Coefficients decouple from scale}

Once we fix the rates $\{r_i\}$, the decay scale $\lambda$, and the gating exponent $q$, the linearized error~\eqref{eq:linearize} is a \emph{linear} function of the coefficients $c_i$. The basis functions $\psi_i(h) = h^q e^{-r_i\lambda(1-h)}$ are known, and the target is $-\delta(h)$ (also known, since it comes from the tower base). So finding the best coefficients is just constrained linear least squares-solvable instantly by a matrix equation.

This is the VARPRO idea: any optimization problem where some parameters enter linearly can be ``projected out,'' leaving only the nonlinear parameters to search over. Here, the coefficients $c_i$ are linear; the nonlinear parameters are just $\lambda$ and $q$. The result:
\begin{equation}
\min_{\lambda,q}\; E_{\mathrm{opt}}(\lambda,q), \qquad E_{\mathrm{opt}}(\lambda,q) = \min_{c:\,\sum c_i = T_k}\;\max_h\;\left|\delta(h)+\textstyle\sum c_i\,\psi_i(h;\lambda,q)\right|.
\end{equation}

For $N = 16$ (the R6 champion), this collapses an 18-dimensional search to a 2D grid over $(\lambda,q)$ with a linear solve at each grid point. On a GPU, thousands of $(\lambda,q)$ candidates are evaluated per kernel launch, making it feasible to screen hundreds of thousands of rate combinations exhaustively.

%==========================================================================
\section{Fully Expanded Formulas}\label{sec:numerical}
%==========================================================================

\noindent\emph{All results verified at 150 decimal digits using \texttt{mpmath} with exact rational Pad\'e coefficients. This section presents the formulas featured in our interactive web demo, from classical to state-of-the-art, in fully expanded form. The complete catalog of all 55 discovered formulas is in Appendix~\ref{app:complete-index}.}

\medskip
\noindent In all formulas below, $a \ge b > 0$ are the semi-axes and
\[
h = \left(\frac{a-b}{a+b}\right)^{\!2}, \qquad x = 1 - h.
\]


%- Ramanujan I -
\noindent\textbf{1. Ramanujan I (1914) - 5{,}400~ppm.}
\begin{equation}
P \approx \pi\!\left(3(a+b) - \sqrt{(3a+b)(a+3b)}\,\right)
\end{equation}

\bigskip

%- Ramanujan II -
\noindent\textbf{2. Ramanujan II (1914) - 402~ppm.}
\begin{equation}
P \approx \pi(a{+}b)\!\left(1 + \frac{3h}{10 + \sqrt{4 - 3h}}\right)
\end{equation}

\bigskip

%- Cantrell -
\noindent\textbf{3. Cantrell (2001) - 412~ppm.}
\begin{equation}
P \approx \pi(a{+}b)\!\left(1 + \frac{3h}{10 + \sqrt{4 - 3h} + (4/\pi - 14/11)\,h^{3/2}}\right)
\end{equation}

\bigskip

%- Ayala-Rendón -
\noindent\textbf{4. Ayala-Raggi \& Rend\'on Mar\'in RII/2exp (2025) - 0.574~ppm.}
\begin{equation}
P = \frac{P_{\text{RII}}}{1 - 3.375{\times}10^{-4}\,e^{-10.297\,x} - 6.481{\times}10^{-5}\,e^{-40.890\,x}}
\end{equation}

\bigskip

%- RII/2exp beater -
\noindent\textbf{5. Mamoun RII/2exp (this work) - 0.530~ppm, 2 exponentials.}
\begin{equation}
P = P_{\text{RII}}\!\left(1 + h^{0.01}\!\left[6.588{\times}10^{-5}\,e^{-40.127\,x} + 3.365{\times}10^{-4}\,e^{-10.265\,x}\right]\right)
\end{equation}

\bigskip

%- RII/3exp -
\noindent\textbf{6. Mamoun RII/3exp - the Mamoun Equation (this work) - 0.083~ppm, 3 exponentials.}
\begin{equation}
P = \frac{\pi(a{+}b)\!\left(1 + \dfrac{3h}{10+\sqrt{4-3h}}\right)}{1 \;-\; h^{7/3}\!\left[\begin{aligned}
&\; 2.858{\times}10^{-4}\,e^{-6.895\,x} \\[-2pt]
&+ 1.065{\times}10^{-4}\,e^{-22.42\,x} \\[-2pt]
&+ 1.031{\times}10^{-5}\,e^{-103.4\,x}
\end{aligned}\right]}
\end{equation}
\begin{center}\small
\begin{tabular}{lll}
$q = 7/3 = a_1/a_0$ & $\lambda = 6.895$ & $T = 1 - 7\pi/22 \approx 4.025 \times 10^{-4}$ \\
Rates: $1,\; 81/25,\; 15$ & & $A+B+C = T$
\end{tabular}
\end{center}

\bigskip

\bigskip
\noindent\textbf{Formulas 7-10: The Convergent Tower champions.}
The Mamoun Equation above uses the RII base (Ramanujan~II itself). By stacking rational Pad\'e corrections-each one a cubic-over-quadratic rational function multiplied by a high power of~$h$-we build successively better base approximations $R_3, R_4, R_5, R_6$. Each level matches 5~more Taylor coefficients of the exact elliptic integral. Combined with more exponential terms, accuracy improves by orders of magnitude.

Below are the four tower-level champions. For each one, we first write out what $R_k(h)$ actually is (the ``rational form''), then give the complete perimeter formula with every exponential term spelled out (the ``coefficient form''). In the coefficient form, $\alpha_i = r_i \cdot \lambda$ has been pre-multiplied so the formula is entirely self-contained. All 55 discovered formulas-from 4 to 20 exponentials across four tower levels-are listed in Appendix~\ref{app:complete-index}.

\bigskip

%- R3+15exp -
\noindent\textbf{7. Mamoun R3+15exp (this work) - 1.427~ppb, 15 exponentials.}

\smallskip\noindent The $R_3$ base is Ramanujan~II plus one Pad\'e correction, matching $B(h)$ through $h^9$.  At $h=1$ this locks on to the convergent $355/113$ of~$\pi$; the exponential correction must use up the error budget $T_3 = 1 - 113\pi/355 = +8.49 \times 10^{-8}$.  The complete formula ($x = 1-h$):
\noindent\emph{Rational base:}
\[
R_3(h) = 1 + \frac{3h}{10+\sqrt{4{-}3h}}
  \;+\; h^5\!\cdot\!\frac{2.289{\times}10^{-5} + 6.604{\times}10^{-6}\,h + 1.959{\times}10^{-6}\,h^2 + 5.209{\times}10^{-7}\,h^3}
       {1 - 1.357\,h + 0.4197\,h^2}
\]
\emph{Full perimeter:}
\begin{equation*}
P = \frac{\pi(a{+}b)\,R_3(h)}{1 - h^{3.416}\!\left(\!\begin{aligned}
  &\; +1.069{\times}10^{-6}\,e^{-6.2414\,x} \; -5.345{\times}10^{-6}\,e^{-9.7089\,x} \\[-3pt]
  &\quad -4.028{\times}10^{-6}\,e^{-14.563\,x} \; -5.070{\times}10^{-6}\,e^{-16.991\,x} \\[-3pt]
  &\quad -5.344{\times}10^{-6}\,e^{-21.845\,x} \; +3.808{\times}10^{-6}\,e^{-30.874\,x} \\[-3pt]
  &\quad +9.540{\times}10^{-6}\,e^{-47.185\,x} \; -3.477{\times}10^{-6}\,e^{-94.370\,x} \\[-3pt]
  &\quad +1.034{\times}10^{-5}\,e^{-110.10\,x} \; -3.937{\times}10^{-6}\,e^{-189.32\,x} \\[-3pt]
  &\quad -3.353{\times}10^{-6}\,e^{-218.45\,x} \; +1.059{\times}10^{-5}\,e^{-283.50\,x} \\[-3pt]
  &\quad -3.064{\times}10^{-6}\,e^{-349.52\,x} \; -5.329{\times}10^{-6}\,e^{-436.90\,x} \\[-3pt]
  &\quad +3.678{\times}10^{-6}\,e^{-509.72\,x}
\end{aligned}\!\right)}, \quad x = 1{-}h
\end{equation*}

\bigskip

%- R4+20exp -
\noindent\textbf{8. Mamoun R4+20exp (this work) - 0.492~ppb, 20 exponentials.}

\smallskip\noindent The $R_4$ base adds a second Pad\'e layer, matching $B(h)$ through $h^{14}$.  At $h=1$ this locks on to the convergent $103993/33102$; the error budget is $T_4 = -1.84 \times 10^{-10}$ (negative because this convergent overshoots $\pi$).

\smallskip\noindent\emph{Rational base:}
\[
R_4(h) = 1 + \frac{3h}{10+\sqrt{4{-}3h}}
  \;+\; h^5\!\cdot\!\frac{2.289{\times}10^{-5} + 6.604{\times}10^{-6}\,h + 1.959{\times}10^{-6}\,h^2 + 5.209{\times}10^{-7}\,h^3}
       {1 - 1.357\,h + 0.4197\,h^2}
\]
\vspace{-8pt}
\[
  \phantom{R_4(h) = {}}
  +\; h^{10}\!\cdot\!\frac{-1.979{\times}10^{-7} + 1.864{\times}10^{-9}\,h + 1.238{\times}10^{-7}\,h^2 + 1.083{\times}10^{-7}\,h^3}
       {1 - 2.680\,h + 2.014\,h^2}
\]
\emph{Full perimeter ($x = 1-h$):}
\begin{equation*}
P = \frac{\pi(a{+}b)\,R_4(h)}{1 - h^{2.520}\!\left(\!\begin{aligned}
  &\; -1.594{\times}10^{-6}\,e^{-5.2674\,x} \; +4.005{\times}10^{-6}\,e^{-6.7724\,x} \\[-3pt]
  &\quad -6.078{\times}10^{-6}\,e^{-11.852\,x} \; -1.789{\times}10^{-5}\,e^{-15.802\,x} \\[-3pt]
  &\quad +7.005{\times}10^{-6}\,e^{-18.436\,x} \; -1.534{\times}10^{-6}\,e^{-36.872\,x} \\[-3pt]
  &\quad +8.604{\times}10^{-6}\,e^{-47.407\,x} \; -5.919{\times}10^{-7}\,e^{-51.199\,x} \\[-3pt]
  &\quad +4.016{\times}10^{-6}\,e^{-73.242\,x} \; +2.309{\times}10^{-6}\,e^{-79.012\,x} \\[-3pt]
  &\quad -4.071{\times}10^{-6}\,e^{-119.47\,x} \; +6.387{\times}10^{-6}\,e^{-142.22\,x} \\[-3pt]
  &\quad -4.451{\times}10^{-7}\,e^{-189.63\,x} \; -7.202{\times}10^{-6}\,e^{-237.03\,x} \\[-3pt]
  &\quad +9.244{\times}10^{-6}\,e^{-268.64\,x} \; -3.986{\times}10^{-6}\,e^{-395.06\,x} \\[-3pt]
  &\quad -1.640{\times}10^{-6}\,e^{-474.07\,x} \; +5.453{\times}10^{-6}\,e^{-553.08\,x} \\[-3pt]
  &\quad -2.320{\times}10^{-6}\,e^{-711.10\,x} \; +3.261{\times}10^{-7}\,e^{-1106.2\,x}
\end{aligned}\!\right)}
\end{equation*}

\bigskip

%=== SCOR COMPACT CHAMPIONS ===
\noindent\textbf{Formulas 9--12: SCOR Compact Champions.}

\smallskip\noindent
SCOR (Singularity-Conscious Optimal Reduction) enriches the correction basis with $x^k\ln x$ terms that absorb the logarithmic singularity of $E(m)$ at $m=1$.  These terms have \textbf{zero free parameters}---their coefficients are solved analytically by VARPRO alongside the exponential coefficients.  The result: dramatically fewer exponentials suffice for the same accuracy.

\medskip\noindent
\textbf{Optimality status.}  The amplitudes (exponential coefficients $c_i$ and log coefficients $d_j$) are \emph{provably} optimal: they solve a convex, linearly constrained least-squares problem with a unique global minimum.  The rates $\mu_i$ and global parameters $\lambda, q$ are found by CMA-ES (a stochastic search) and are \emph{numerically} near-optimal (verified by multi-restart convergence) but not provably global.

\bigskip

%--- Formula 9 ---
\noindent\textbf{9. Mamoun $R_3$+2exp+1log (this work) --- 34~ppb, 4 free parameters.}

\smallskip\noindent Same parameter count as Ayala-Rend\'on (2025), but $17\times$ better accuracy.  Fully self-contained ($x = 1-h$):
\begin{equation*}
\boxed{P = \frac{\pi(a{+}b)\!\left(1 + \dfrac{3h}{10+\sqrt{4{-}3h}}
+ h^5\!\cdot\!\dfrac{2.289{\times}10^{-5} + 6.604{\times}10^{-6}h + 1.959{\times}10^{-6}h^2 + 5.209{\times}10^{-7}h^3}
{1 - 1.357h + 0.4197h^2}\right)}
{1 - h^{18.947}\!\left(
-7.638{\times}10^{-6}\,e^{-11.952\,x}
+7.638{\times}10^{-6}\,e^{-78.315\,x}
+7.163{\times}10^{-5}\,x\ln x
\right)}}
\end{equation*}

\bigskip

%--- Formula 10 ---
\noindent\textbf{10. Mamoun $R_5$+3exp+3log (this work) --- 2.3~ppb, 5 free parameters.}

\smallskip\noindent Uses the full $R_5$ tower (Ramanujan~II + 3 Pad\'e layers).  Fully expanded:
\[
R_5(h) = 1 + \frac{3h}{10+\sqrt{4{-}3h}}
  + h^5\!\cdot\!\frac{2.289{\times}10^{-5} + 6.604{\times}10^{-6}h + 1.959{\times}10^{-6}h^2 + 5.209{\times}10^{-7}h^3}
       {1 - 1.357h + 0.4197h^2}
\]
\vspace{-8pt}
\[
  + h^{10}\!\cdot\!\frac{-1.979{\times}10^{-7} + 1.864{\times}10^{-9}h + 1.238{\times}10^{-7}h^2 + 1.083{\times}10^{-7}h^3}
       {1 - 2.680h + 2.014h^2}
  + h^{15}\!\cdot\!\frac{-1.472{\times}10^{-7} - 1.167{\times}10^{-7}h + 5.780{\times}10^{-8}h^2 + 2.056{\times}10^{-7}h^3}
       {1 - 3.886h + 4.420h^2}
\]
\emph{Full perimeter:}
\begin{equation*}
\boxed{P = \frac{\pi(a{+}b)\,R_5(h)}{1 - h^{21.961}\!\left(\!\begin{aligned}
  &-2.972{\times}10^{-6}\,e^{-15.350\,x}
  +3.041{\times}10^{-6}\,e^{-101.64\,x}
  -6.834{\times}10^{-8}\,e^{-4032.6\,x} \\[2pt]
  &+1.112{\times}10^{-4}\,x\ln x
  -5.213{\times}10^{-4}\,x^2\ln x
  +3.300{\times}10^{-3}\,x^3\ln x
\end{aligned}\!\right)}}
\end{equation*}

\bigskip

%--- Formula 11 ---
\noindent\textbf{11. Mamoun $R_3$+4exp+3log (this work) --- 0.67~ppb, 6 free parameters.}

\smallskip\noindent Sub-ppb with only 6 free parameters.  Same $R_3$ base as Formula~9:
\begin{equation*}
\boxed{P = \frac{\pi(a{+}b)\,R_3(h)}{1 - h^{20.393}\!\left(\!\begin{aligned}
  &-6.696{\times}10^{-6}\,e^{-0.362\,x}
  +5.756{\times}10^{-6}\,e^{-63.52\,x}
  +9.676{\times}10^{-7}\,e^{-200.9\,x}
  -2.751{\times}10^{-8}\,e^{-5949\,x} \\[2pt]
  &+7.077{\times}10^{-5}\,x\ln x
  -2.041{\times}10^{-4}\,x^2\ln x
  +1.563{\times}10^{-3}\,x^3\ln x
\end{aligned}\!\right)}}
\end{equation*}

\bigskip

%--- Formula 12: GRAND CHAMPION ---
\noindent\textbf{12. Mamoun $R_6$+16exp+3log --- GRAND CHAMPION (this work) --- 0.007~ppb.}

\smallskip\noindent The full $R_6$ tower (4 Pad\'e layers, 25 Taylor coefficients matched) with 16 exponentials and 3 logarithmic enrichment terms.  Verified at 300 decimal digits.  The gate exponent $q = 12.4$ suppresses the correction at moderate eccentricity, concentrating all capacity on the extreme-eccentricity tail where the tower error is largest.

\smallskip\noindent\emph{Rational base --- every term written out, nothing to ``plug in'':}

Start with Ramanujan's formula:
\[
R_2(h) = 1 + \frac{3h}{10 + \sqrt{4 - 3h}}
\]
Then add four correction layers, each a fraction multiplied by a high power of~$h$:
\begin{align*}
R_3(h) &= R_2(h) \;+\; h^5 \times
  \frac{0.00002289 + 0.000006604\,h + 0.000001959\,h^2 + 0.0000005209\,h^3}
       {1 - 1.3573\,h + 0.4197\,h^2}
\\[8pt]
R_4(h) &= R_3(h) \;+\; h^{10} \times
  \frac{-1.979{\times}10^{-7} + 1.864{\times}10^{-9}\,h + 1.238{\times}10^{-7}\,h^2 + 1.083{\times}10^{-7}\,h^3}
       {1 - 2.6800\,h + 2.0136\,h^2}
\\[8pt]
R_5(h) &= R_4(h) \;+\; h^{15} \times
  \frac{-1.472{\times}10^{-7} - 1.167{\times}10^{-7}\,h + 5.780{\times}10^{-8}\,h^2 + 2.056{\times}10^{-7}\,h^3}
       {1 - 3.8862\,h + 4.4201\,h^2}
\\[8pt]
R_6(h) &= R_5(h) \;+\; h^{20} \times
  \frac{-1.057{\times}10^{-6} - 1.828{\times}10^{-6}\,h - 4.965{\times}10^{-7}\,h^2 + 3.382{\times}10^{-6}\,h^3}
       {1 - 5.5893\,h + 9.3243\,h^2}
\end{align*}
\emph{Note:} Each correction is multiplied by $h^{5k}$ (for $k = 1, 2, 3, 4$), so near $h = 0$ (circles) all corrections vanish and $R_6 \approx R_2$ (Ramanujan).  Near $h = 1$ (flat ellipses) all corrections contribute.  The formula is purely algebraic: only addition, multiplication, division, and one square root.

\emph{Full perimeter ($x = 1-h$):}
\begin{equation}\label{eq:grand-champion}
\boxed{P = \frac{\pi(a{+}b)\,R_6(h)}{1 - h^{12.415}\!\left(\!\begin{aligned}
  &\; +2.662{\times}10^{-4}\,e^{-\phantom{0}3.185\,x}
     -8.682{\times}10^{-4}\,e^{-12.437\,x}
     +1.417{\times}10^{-3}\,e^{-21.152\,x} \\[-3pt]
  &\; -1.523{\times}10^{-3}\,e^{-23.906\,x}
     +1.050{\times}10^{-3}\,e^{-35.041\,x}
     -5.408{\times}10^{-4}\,e^{-42.430\,x} \\[-3pt]
  &\; +1.930{\times}10^{-4}\,e^{-79.284\,x}
     +3.535{\times}10^{-6}\,e^{-329.32\,x}
     +1.156{\times}10^{-6}\,e^{-651.07\,x} \\[-3pt]
  &\; +6.182{\times}10^{-7}\,e^{-1163.3\,x}
     -3.132{\times}10^{-7}\,e^{-2090.2\,x}
     +3.990{\times}10^{-7}\,e^{-2503.1\,x} \\[-3pt]
  &\; -2.266{\times}10^{-7}\,e^{-4104.5\,x}
     +1.279{\times}10^{-7}\,e^{-5109.6\,x}
     -8.260{\times}10^{-9}\,e^{-17090\,x} \\[-3pt]
  &\; +1.099{\times}10^{-8}\,e^{-18340\,x} \\[3pt]
  &\; +1.902{\times}10^{-6}\,x\ln x
     -4.374{\times}10^{-4}\,x^2\ln x
     +1.178{\times}10^{-3}\,x^3\ln x
\end{aligned}\!\right)}}
\end{equation}

\bigskip

\noindent\textbf{Complete leaderboard:}
\begin{center}
\begin{tabular}{@{}rlcccrc@{}}
\toprule
\# & Formula & Tower & Exp & Log & Free & Error \\
\midrule
  & Ramanujan~II (1914)       & $R_2$ & 0  & 0 &  0 & 402~ppm \\
  & Ayala-Rend\'on (2025)     & $R_2$ & 2  & 0 &  4 & 0.574~ppm \\
6 & Mamoun Equation           & $R_2$ & 3  & 0 &  1 & 0.083~ppm \\
9 & $R_3$+2exp+1log           & $R_3$ & 2  & 1 &  4 & 34~ppb \\
10& $R_5$+3exp+3log           & $R_5$ & 3  & 3 &  5 & 2.3~ppb \\
11& $R_3$+4exp+3log           & $R_3$ & 4  & 3 &  6 & 0.67~ppb \\
\textbf{12}& $R_6$+16exp+3log  & $R_6$ & 16 & 3 & 18 & \textbf{0.007~ppb} \\
\bottomrule
\end{tabular}
\end{center}
\noindent \emph{Free} = optimization parameters (rates + $\lambda$ + $q$).
Tower coefficients (determined by Taylor matching) and log coefficients (determined by Frobenius expansion) cost zero search effort.

%==========================================================================
\section{Open Questions}\label{sec:open}
%==========================================================================

\textbf{Problem 1.} \emph{Characterize the Franel crossover.} The deficit ratio $\delta_{n+1}/\delta_n$ passes through $7/8$ near $n = 20$ (Theorem~\ref{thm:deficit}) before converging to~$1$.  Give a closed-form expression for the crossover index as a function of the CF quotients.

\begin{conjecture}[Tower convergence]
$R_k(1) \to 4/\pi$ as $k \to \infty$: the tower exhausts all convergents of~$\pi$.
\end{conjecture}

\begin{conjecture}[Optimal exponential convergence rate]
The minimax error of $R_k + N$-exp scales as $C_k \cdot N^{-\alpha_k}$
with $\alpha_k$ independent of~$k$ (polynomial, not exponential, in~$N$),
and $C_k \sim |T_k|$.
\end{conjecture}

\textbf{Problem 4.} \emph{Give a rigorous lower bound} on the minimax error of any $N$-exponential correction to $R_k$.

\textbf{Problem 5.} \emph{Explain pair-wise rate necessity.} Is there an a priori criterion predicting when pair-wise rate search is needed (e.g., from the sign of $T_k$)?

\textbf{Problem 6.} \emph{Does the VARPRO 2D landscape $E_{\mathrm{opt}}(\lambda,q)$ have a unique global minimum?} Numerical evidence suggests yes for each rate set, but a proof would guarantee that the 2D search finds the true optimum.

\textbf{Problem 7.} \emph{Denominator closed form for $18/(-8+\pi^2)$.} Theorem~\ref{thm:cf-closed} gives the numerator $\tilde{p}(n)$ in closed form.  Does the denominator $\tilde{q}(n)$ also admit one?

%==========================================================================
\section{Conclusion}\label{sec:conclusion}
%==========================================================================

We have shown that the ellipse perimeter problem admits a systematic hierarchy of approximations, each level controlled by $\pi$'s continued fraction. The key results are:

\begin{enumerate}
\item Ramanujan~II is a CF formula and the unique $(2,1,1)$ QHP of $B(h)$.
\item The convergent tower adds 5 Taylor coefficients per level via constrained Pad\'e.
\item The exponential correction separates linearly (VARPRO), reducing 18D to 2D.
\item R6+16exp+3log (SCOR) achieves $0.007$~ppb---the best verified minimax error, far \emph{below} the na\"{\i}ve BARF (which is a pointwise, not minimax, bound for purely analytic corrections).
\item The BARF and the Focal Uncertainty Principle set structural limits; the gate exponent $q$ provides an additional degree of freedom that circumvents the pointwise BARF by redirecting approximation effort away from where the residue peaks.
\item A closed form for the Euler-Wallis numerators of $18/(-8+\pi^2)$ links central binomial coefficients to~$\pi^2$.
\end{enumerate}

\noindent\textbf{Data availability.} Complete parameter sets for all results are available at \url{https://github.com/monir/ellipse} in machine-readable CSV format.

%==========================================================================
\section*{Acknowledgments}
%==========================================================================

Dedicated to my wife Sarah, who graciously put up with weeks of my staring at the walls and muttering incoherently during the course of this research.

\medskip
Computations performed with \texttt{mpmath} (150 digits), CMA-ES, PyTorch (5$\times$ RTX 4090 GPU cluster), and Python \texttt{Fraction} (exact rational arithmetic).

%==========================================================================
\appendix
\section{Computational Methodology and Reproducibility}\label{app:reproducibility}
%==========================================================================

This appendix documents the computational pipeline used to obtain all numerical results, organized by development stage. All code is available at \url{https://github.com/monir/ellipse}.

\subsection{Stage 1: Tower construction}

The Pad\'e tower (Theorem~\ref{thm:tower}) is built by \texttt{tower\_gpu\_kernel.py}, function \texttt{build\_tower(k)}. For level $k$, it computes the constrained $[3,2]$ Pad\'e approximant using Python \texttt{Fraction} arithmetic (exact rationals) and returns the layer coefficients $N_i$, $D_i$, onset exponents, and the correction budget $T_k$.

\medskip\noindent\textbf{To reproduce:}
\begin{verbatim}
from tower_gpu_kernel import build_tower
layers = build_tower(5)   # level 5 -> R6
T_k = layers[-1]['T']     # T_k ~ -3.89e-11
\end{verbatim}

\subsection{Stage 2: Rate discovery via CMA-ES}

The initial optimization used CMA-ES (Covariance Matrix Adaptation Evolution Strategy) implemented entirely in PyTorch on GPU. Each candidate rate set $\{\mu_1,\ldots,\mu_N\}$ defines an $(N{+}2)$-dimensional search over $N$ ratio parameters, $\lambda$, and $q$.

The GPU kernel (\texttt{tower\_gpu\_kernel\_v7.py}, class \texttt{TowerGPUKernelV7}) evaluates $\max|\varepsilon(h)|$ over a uniform grid of $G$ points in $h \in (0,1)$, all in FP64 arithmetic. The CMA-ES optimizer (\texttt{GPUCMAES} class) runs entirely on-device with no CPU round-trips during optimization.

\medskip\noindent\textbf{Key scripts (in chronological development order):}
\begin{itemize}
\item \texttt{rk\_tower\_quad\_v10.py}: Discovered the R6+16-exp configuration ($0.409$~ppb). Tested all $\binom{49}{4} = 211{,}876$ four-rate additions to the 12-rate R6 base. Runtime: 4.7~hours on local hardware.
\item \texttt{rk\_warmstart\_campaign\_v5.py}: Mac MPS warm-start campaign that explored R6+18-exp ($0.358$~ppb on coarse grid). Add-2 screen (C3 phase) found rates $\{18, 50\}$. True peak error 115~ppb at extreme eccentricity---not a champion.
\item \texttt{rk\_warmstart\_campaign\_v1.py}: Multi-phase campaign with warm-starting. Phase C1 re-polished the existing champion rates with fresh CMA-ES restarts, discovering improved coefficients at $0.015$~ppb (same 16 rates).
\item \texttt{rk\_warmstart\_campaign\_v2.py}: Extended campaign with Franel 4-block structure (one rate per zone: slow/medium/fast/ultra) for systematic exploration.
\item \texttt{r6\_hierarchical\_optimizer\_v1.py}: Four-phase hierarchical optimizer (VARPRO + L-BFGS-B + coordinate golden-section + CMA-ES local polish) that treats all 16 rates as continuous variables and jointly optimises $(\mu_1,\ldots,\mu_{16},\lambda,q)$. Improved the pure-exponential R6+16exp result from $0.058$~ppb to $0.049$~ppb, later refined to $0.018$~ppb. The overall champion is the SCOR variant at $0.007$~ppb (Formula~12).
\end{itemize}

\medskip\noindent\textbf{Rate pool.} The candidate rates are drawn from a pool of ${\sim}65$ values motivated by the continued fraction structure of $\pi$: $\{1/3, 3/7, 1/2, 2/3, 3/4, 1, 7/6, \ldots, 200\}$. CF-derived rates include $81/25 = a_0^4/(a_0{+}a_1{+}a_2)$, $292/63$, and $45/7$.

\medskip\noindent\textbf{Hardware.} Linux server with 5$\times$ NVIDIA RTX 4090 (24~GB VRAM each), 256~GB system RAM, CUDA 12.8, PyTorch 2.10. Mac Studio M2 Ultra (192~GB) and MacBook Pro M2 Max (96~GB) for verification, symbolic work, and local LLM inference. All three machines ran computational loads throughout the research. Typical GPU throughput: 5.7 rate combinations per second at population 512, 20K-point evaluation grid.

\subsection{Stage 3: VARPRO screening}

The VARPRO separation (\S\ref{sec:varpro}) was implemented as a GPU-accelerated screening pipeline (\texttt{rk\_varpro\_screen\_v2.py}) that eliminates the CMA-ES bottleneck for coefficient optimization.

\medskip\noindent\textbf{Algorithm.} For fixed rates and $(\lambda,q)$, the constrained least-squares problem $\min \|\delta + \Phi c\|_2^2$ subject to $\sum c_i = T_k$ is solved via variable elimination: setting $c_N = T_k - \sum_{i<N} c_i$ reduces to an unconstrained $(N{-}1)$-dimensional normal equation $\tilde{A}\tilde{c} = \tilde{b}$.

\medskip\noindent\textbf{Gram matrix optimization.} Rather than forming the full $(B \times N \times G)$ basis tensor per batch (requiring $G$-dimensional matrix multiplications per combination), we precompute $R \times R$ Gram matrices:
\[
\text{Gram}_{l}[i,j] = \sum_{g=1}^{G} \psi_i(h_g;\lambda_l,q_l)\,\psi_j(h_g;\lambda_l,q_l),
\qquad
\text{proj}_l[i] = \sum_{g=1}^{G} \psi_i(h_g;\lambda_l,q_l)\,\delta(h_g),
\]
where $R \approx 65$ is the total rate pool size. Per-combination work then reduces to extracting an $N \times N$ submatrix and solving a system-$O(N^3)$ with \emph{zero} $G$-dimensional operations.

\medskip\noindent\textbf{Three-phase pipeline:}
\begin{enumerate}
\item \textbf{P1: Gram-based $L^2$ screening.} All $\binom{49}{4} = 211{,}876$ combinations screened over $12 \times 16 = 192$ coarse $(\lambda,q)$ grid points. Runtime: $\sim$8~seconds on a single GPU. Throughput: ${\sim}25{,}000$ combinations/second.
\item \textbf{P2: $L^\infty$ refinement.} Top 200-500 combinations evaluated with full $G$-dimensional residual computation on a dense $44 \times 31 = 1{,}364$-point $(\lambda,q)$ grid. Runtime: $\sim$3~minutes.
\item \textbf{P3: CMA-ES polish.} Top 50 combinations polished with full $(N{+}2)$-dimensional CMA-ES (population 4096, 200K function evaluations, 3 restarts each) on a 20K-point evaluation grid. Runtime: $\sim$1~minute on 5 GPUs.
\end{enumerate}

\noindent Total pipeline: $\sim$5~minutes vs $\sim$10~hours for pure CMA-ES screening (${\sim}120\times$ speedup).

\subsection{Stage 4: Verification}

All champion results are verified independently:
\begin{enumerate}
\item \textbf{150-digit verification.} Using \texttt{mpmath} with 150 decimal digits, the exact perimeter $P = 4a\,E(e^2)$ is computed at 5{,}000 uniformly-spaced eccentricity values and compared with the formula output. The reported $\max|\varepsilon|$ is the supremum over this grid.
\item \textbf{Exact rational tower coefficients.} All Pad\'e tower layers use Python \texttt{Fraction} arithmetic (arbitrary-precision rationals), eliminating floating-point error in the tower construction. The exponential correction parameters are the only floating-point quantities.
\item \textbf{Cross-validation.} The R6 base error profile $\delta(h) = P_{R_6}(h)/P_{\text{exact}}(h) - 1$ is independently computed by \texttt{formulas.py} and the GPU kernel, with agreement to 15 significant digits.
\end{enumerate}

\subsection{File manifest}

\newcommand{\fp}[1]{\texttt{#1}}% file path shorthand
\begin{center}\small
\begin{longtable}{@{}l >{\raggedright\arraybackslash}p{7.8cm}@{}}
\toprule
File & Purpose \\
\midrule
\endfirsthead
\toprule
File & Purpose \\
\midrule
\endhead
\fp{src/formulas.py}               & All perimeter formulas (Ramanujan~I/II, tower $R_2$--$R_6$, exponential corrections) \\[3pt]
\fp{demo/ellipse\_web\_demo.py}     & Interactive web demo: comparison charts, winner-band visualization, error analysis \\[3pt]
\fp{gpu\_.../tower\_gpu\_kernel.py}  & Tower construction (\fp{build\_tower}) and base GPU evaluation \\[3pt]
\fp{gpu\_.../tower\_gpu\_kernel\_v7.py} & Production GPU kernel: \fp{TowerGPUKernelV7} (FP64 batch, \fp{GPUCMAES}) \\[3pt]
\fp{r\_scripts/rk\_tower\_quad\_v10.py} & Original quad search that found R6+16-exp ($0.409$~ppb) \\[3pt]
\fp{r\_scripts/rk\_warmstart\_v5.py}    & Mac MPS add-2 campaign; explored R6+18-exp ($0.358$~ppb coarse) \\[3pt]
\fp{r\_scripts/rk\_warmstart\_v1.py}    & 4-phase warm-start campaign (C1--C4) \\[3pt]
\fp{r\_scripts/rk\_varpro\_screen\_v1.py} & VARPRO screening, first implementation (BMM per combo) \\[3pt]
\fp{r\_scripts/rk\_varpro\_screen\_v2.py} & VARPRO screening with Gram matrix optimization \\[3pt]
\fp{r\_scripts/outputs/*.json}          & Machine-readable results for all runs \\[3pt]
\fp{active/best\_results.csv}           & Summary CSV of all champion results \\
\bottomrule
\end{longtable}
\end{center}

\medskip\noindent\textbf{To replicate the R6+16-exp result (pure exponential, 0.018~ppb):}
\begin{verbatim}
# On Apple MPS or CUDA GPU:
cd ellipse/active/research_scripts
python3 -u rk_warmstart_campaign_v5.py --skip-c3  # runs C1+C2+C6 only
# Or for the full add-2 screen (finds the champion):
python3 -u rk_warmstart_campaign_v5.py            # ~22 min on M2 Ultra

# To evaluate just the known champion:
python3 -c "
from gpu_tower_server.tower_gpu_kernel_v7 import *
from gpu_tower_server.tower_gpu_kernel import build_tower
import torch
layers = build_tower(5)
T_k = layers[-1]['T']
k = TowerGPUKernelV7(layers, T_k,
    device='cuda:0', grid_pts=20000)
rates = torch.tensor([1/3,3/7,3/4,1,1.5,3,3.24,7.56,
    15,17,20,30,35,36,45,55],
    dtype=torch.float64, device='cuda:0')
x = torch.tensor([-2.83,9.06,0.47,-0.27,-0.45,
    -5.16,-1.54,2.45,-5.87,4.13,-1.26,
    0.73,-1.33,1.46,-0.60,17.575,2.591],
    dtype=torch.float64, device='cuda:0').unsqueeze(0)
err = k.batch_eval(rates, x, mode='mixed')
print(f'Error: {err.item()*1000:.3f} ppb')
"
\end{verbatim}

%==========================================================================
\begin{thebibliography}{99}
%==========================================================================

\bibitem{ramanujan1914}
S.~Ramanujan, ``Modular equations and approximations to $\pi$,''
\textit{Q.\ J.\ Math.}, vol.~45, pp.~350-372, 1914.

\bibitem{villarino2005}
M.~B.~Villarino, ``Ramanujan's perimeter of an ellipse,''
arXiv:math/0506384, 2005.

\bibitem{almkvist1988}
G.~Almkvist and B.~C.~Berndt, ``Gauss, Landen, Ramanujan, the
arithmetic-geometric mean, ellipses, $\pi$, and the Ladies Diary,''
\textit{Am.\ Math.\ Mon.}, vol.~95, pp.~585-608, 1988.

\bibitem{ayala2025}
S.~E.~Ayala Raggi and M.~Rend\'on Mar\'in, ``New closed form approximations to the
perimeter of an ellipse,'' arXiv:2510.16191, 2025.

\bibitem{golub2003}
G.~H.~Golub and V.~Pereyra, ``Separable nonlinear least squares: the variable projection method and its applications,''
\textit{Inverse Problems}, vol.~19, pp.~R1-R26, 2003.

\bibitem{wang2026}
C.~Wang, ``A formal proof of a continued fraction conjecture for $\pi$ originating from the Ramanujan Machine,''
arXiv:2601.08461, 2026.

\bibitem{newman1964}
D.~J.~Newman, ``Rational approximation to $|x|$,''
\textit{Michigan Mathematical Journal}, vol.~11, pp.~11--14, 1964.

\bibitem{zwdeng2023}
Z.~Wang and Y.~Deng, ``On periodicity of continued fractions with partial
quotients in quadratic number fields,'' arXiv:2304.11803, 2023.

\end{thebibliography}

%==========================================================================
\section{Complete Index of Discovered Formulas}\label{app:complete-index}
%==========================================================================

\subsection*{How to read these formulas}

\noindent\textbf{Variables.} Given an ellipse with semi-major axis~$a$ and semi-minor axis~$b$:
\[
h = \left(\frac{a-b}{a+b}\right)^{\!2} \quad (= 0 \text{ for a circle},\; \to 1 \text{ as } b \to 0), \qquad x = 1-h.
\]

\noindent\textbf{The rational base $R_k(h)$.} Each formula uses a rational function $R_k(h)$ as its starting approximation, built by stacking Pad\'e corrections on top of Ramanujan's second formula. Each correction is a cubic-over-quadratic rational function $N_k(h)/D_k(h)$ with exact rational coefficients (Theorem~\ref{thm:tower}), adding 5 Taylor terms per level.  Here they are explicitly (decimals shown; exact rationals in the code repository):
\begin{align*}
R_2(h) &= 1 + \frac{3h}{10 + \sqrt{4-3h}} & &\text{(Ramanujan II)},
\end{align*}
\begin{align*}
R_3(h) &= R_2(h) + h^5 \cdot \frac{2.2889{\times}10^{-5} + 6.6038{\times}10^{-6}\,h + 1.9585{\times}10^{-6}\,h^2 + 5.2090{\times}10^{-7}\,h^3}{1 - 1.35731\,h + 0.41973\,h^2},
\end{align*}
\begin{align*}
R_4(h) &= R_3(h) + h^{10} \cdot \frac{-1.9786{\times}10^{-7} + 1.8641{\times}10^{-9}\,h + 1.2381{\times}10^{-7}\,h^2 + 1.0833{\times}10^{-7}\,h^3}{1 - 2.68000\,h + 2.01361\,h^2},
\end{align*}
\begin{align*}
R_5(h) &= R_4(h) + h^{15} \cdot \frac{-1.4724{\times}10^{-7} - 1.1670{\times}10^{-7}\,h + 5.7803{\times}10^{-8}\,h^2 + 2.0557{\times}10^{-7}\,h^3}{1 - 3.88616\,h + 4.42007\,h^2},
\end{align*}
\begin{align*}
R_6(h) &= R_5(h) + h^{20} \cdot \frac{-1.0568{\times}10^{-6} - 1.8281{\times}10^{-6}\,h - 4.9648{\times}10^{-7}\,h^2 + 3.3822{\times}10^{-6}\,h^3}{1 - 5.58931\,h + 9.32427\,h^2}.
\end{align*}
Each successive correction is multiplied by $h^{5k}$, so it is negligible near $h=0$ (circles) and only matters near $h=1$ (flat ellipses). The $R_k$ base matches the exact function $B(h) = {}_2F_1(-\tfrac12,-\tfrac12;1;h)$ through order~$h^{5k-6}$.

\medskip
\noindent\textbf{Endpoint constraint.} At $h=1$ (the degenerate ellipse $b=0$), the exact perimeter is $P = 4a$, so $B(1) = 4/\pi$.  Each $R_k$ is constrained so that $R_k(1)$ yields the ratio $4q_j/p_j$ for a specific convergent $p_j/q_j$ of~$\pi$-the tower ``locks on'' to successively better rational approximations of~$\pi$:
\begin{center}\small
\begin{tabular}{llll}
\toprule
Base & $\pi(a+b)\,R_k(1)$ yields & Convergent of $\pi$ & Error budget $T_k$ \\
\midrule
$R_2$ & $4a \times 7/22$ & $22/7$ & $+4.0 \times 10^{-4}$ \\
$R_3$ & $4a \times 113/355$ & $355/113$ & $+8.5 \times 10^{-8}$ \\
$R_4$ & $4a \times 33102/103993$ & $103993/33102$ & $-1.8 \times 10^{-10}$ \\
$R_5$ & $4a \times 33215/104348$ & $104348/33215$ & $+1.1 \times 10^{-10}$ \\
$R_6$ & next convergent & & ${\sim}10^{-11}$ \\
\bottomrule
\end{tabular}
\end{center}
The \emph{error budget} $T_k = 1 - q_j\pi/p_j$ is how far this convergent is from $\pi$; it alternates in sign and shrinks by ${\sim}10^3$ per level. The exponential correction must exactly use up this budget.

\medskip
\noindent\textbf{The exponential correction.} The complete perimeter formula is:
\begin{equation*}
\boxed{P = \frac{\pi(a+b)\,R_k(h)}{1 - h^q \displaystyle\sum_{i=1}^{N} c_i\,e^{-\alpha_i\,x}},
\qquad x = 1-h, \quad \alpha_i = r_i \cdot \lambda}
\end{equation*}
where:
\begin{itemize}
\item $q$ is the \textbf{gating exponent}: the factor $h^q$ suppresses the correction near $h = 0$ (circles, where Ramanujan is already excellent) and activates it near $h = 1$ (extreme eccentricities). Larger $q$ means the correction ``turns on'' later.
\item $\lambda$ is the \textbf{global decay scale}: controls how fast each exponential falls off. Larger $\lambda$ means sharper, more localized corrections.
\item $r_i$ are the \textbf{rates}: rational numbers (mostly from the continued fraction of $\pi$) that set the relative decay speeds.  The product $\alpha_i = r_i \cdot \lambda$ is the actual exponent coefficient in each term.
\item $c_i$ are the \textbf{coefficients}: determined by constrained least-squares once $\lambda$, $q$, and the rates are fixed, subject to $\sum c_i = T_k$ (the error budget).
\end{itemize}
In the equations below, $\alpha_i = r_i \cdot \lambda$ is pre-multiplied so each formula is self-contained-you can evaluate it directly without knowing $r_i$ and $\lambda$ separately. Champions (best per tower level) are in \textbf{bold}.

\small

\medskip\noindent\textbf{R3 family} ($R_3$ base, matching $B(h)$ through $h^9$, error budget $T_3 = +8.49 \times 10^{-8}$)\nopagebreak

\noindent In every R3 formula below, $R_3(h)$ stands for:
\[
R_3(h) = 1 + \frac{3h}{10+\sqrt{4{-}3h}}
  \;+\; h^5\!\cdot\!\frac{2.289{\times}10^{-5} + 6.604{\times}10^{-6}\,h + 1.959{\times}10^{-6}\,h^2 + 5.209{\times}10^{-7}\,h^3}
       {1 - 1.357\,h + 0.4197\,h^2}
\]

\noindent\textbf{R3+4} 62.147~ppb\nopagebreak
\begin{equation*}
P = \frac{\pi(a{+}b)\,R_3(h)}{1 - h^{9.9200}\!\left(\!\begin{aligned}
  &\; +2.847{\times}10^{-6}\,e^{-3.3157\,x} \; -2.277{\times}10^{-5}\,e^{-7.7365\,x} \; +1.303{\times}10^{-5}\,e^{-25.066\,x} \\[-3pt]
  &\quad +6.978{\times}10^{-6}\,e^{-116.05\,x}
\end{aligned}\!\right)}, \quad x = 1{-}h
\end{equation*}

\noindent\textbf{R3+5} 59.984~ppb\nopagebreak
\begin{equation*}
P = \frac{\pi(a{+}b)\,R_3(h)}{1 - h^{9.7942}\!\left(\!\begin{aligned}
  &\; +1.369{\times}10^{-6}\,e^{-3.3786\,x} \; -1.095{\times}10^{-5}\,e^{-7.8833\,x} \; -1.095{\times}10^{-5}\,e^{-9.1972\,x} \\[-3pt]
  &\quad +1.384{\times}10^{-5}\,e^{-25.542\,x} \; +6.771{\times}10^{-6}\,e^{-118.25\,x}
\end{aligned}\!\right)}, \quad x = 1{-}h
\end{equation*}

\noindent\textbf{R3+6} 41.470~ppb\nopagebreak
\begin{equation*}
P = \frac{\pi(a{+}b)\,R_3(h)}{1 - h^{3.7378}\!\left(\!\begin{aligned}
  &\; +1.084{\times}10^{-6}\,e^{-5.2127\,x} \; -8.669{\times}10^{-6}\,e^{-12.163\,x} \; -8.664{\times}10^{-6}\,e^{-14.190\,x} \\[-3pt]
  &\quad +9.366{\times}10^{-6}\,e^{-39.408\,x} \; +4.689{\times}10^{-6}\,e^{-91.952\,x} \; +2.279{\times}10^{-6}\,e^{-182.44\,x}
\end{aligned}\!\right)}, \quad x = 1{-}h
\end{equation*}

\noindent\textbf{R3+7} 12.477~ppb\nopagebreak
\begin{equation*}
P = \frac{\pi(a{+}b)\,R_3(h)}{1 - h^{5.5097}\!\left(\!\begin{aligned}
  &\; +1.010{\times}10^{-6}\,e^{-4.2223\,x} \; -8.079{\times}10^{-6}\,e^{-9.8521\,x} \; -7.379{\times}10^{-6}\,e^{-11.494\,x} \\[-3pt]
  &\quad +2.680{\times}10^{-6}\,e^{-31.921\,x} \; +1.212{\times}10^{-5}\,e^{-74.482\,x} \; -3.065{\times}10^{-6}\,e^{-147.78\,x} \\[-3pt]
  &\quad +2.796{\times}10^{-6}\,e^{-236.45\,x}
\end{aligned}\!\right)}, \quad x = 1{-}h
\end{equation*}

\noindent\textbf{R3+8} 4.943~ppb\nopagebreak
\begin{equation*}
P = \frac{\pi(a{+}b)\,R_3(h)}{1 - h^{3.9512}\!\left(\!\begin{aligned}
  &\; +1.415{\times}10^{-6}\,e^{-7.4613\,x} \; -1.131{\times}10^{-5}\,e^{-11.606\,x} \; -1.126{\times}10^{-5}\,e^{-17.410\,x} \\[-3pt]
  &\quad +7.100{\times}10^{-6}\,e^{-20.311\,x} \; +1.051{\times}10^{-5}\,e^{-56.408\,x} \; +3.305{\times}10^{-6}\,e^{-131.62\,x} \\[-3pt]
  &\quad -4.230{\times}10^{-7}\,e^{-261.15\,x} \; +7.466{\times}10^{-7}\,e^{-417.83\,x}
\end{aligned}\!\right)}, \quad x = 1{-}h
\end{equation*}

\noindent\textbf{R3+9} 2.707~ppb\nopagebreak
\begin{equation*}
P = \frac{\pi(a{+}b)\,R_3(h)}{1 - h^{3.0434}\!\left(\!\begin{aligned}
  &\; +1.158{\times}10^{-6}\,e^{-6.6313\,x} \; -5.789{\times}10^{-6}\,e^{-10.315\,x} \; -5.741{\times}10^{-6}\,e^{-15.473\,x} \\[-3pt]
  &\quad -5.521{\times}10^{-6}\,e^{-18.052\,x} \; +1.154{\times}10^{-5}\,e^{-50.133\,x} \; +3.303{\times}10^{-6}\,e^{-116.98\,x} \\[-3pt]
  &\quad +1.571{\times}10^{-6}\,e^{-232.10\,x} \; -1.906{\times}10^{-6}\,e^{-371.35\,x} \; +1.475{\times}10^{-6}\,e^{-464.19\,x}
\end{aligned}\!\right)}, \quad x = 1{-}h
\end{equation*}

\noindent\textbf{R3+10} 2.434~ppb\nopagebreak
\begin{equation*}
P = \frac{\pi(a{+}b)\,R_3(h)}{1 - h^{3.5054}\!\left(\!\begin{aligned}
  &\; +1.157{\times}10^{-6}\,e^{-6.3121\,x} \; -5.493{\times}10^{-6}\,e^{-9.8189\,x} \; -5.747{\times}10^{-6}\,e^{-14.728\,x} \\[-3pt]
  &\quad -5.662{\times}10^{-6}\,e^{-17.183\,x} \; +9.683{\times}10^{-6}\,e^{-47.720\,x} \; +9.314{\times}10^{-6}\,e^{-95.440\,x} \\[-3pt]
  &\quad -5.786{\times}10^{-6}\,e^{-111.35\,x} \; +4.133{\times}10^{-6}\,e^{-220.92\,x} \; -3.902{\times}10^{-6}\,e^{-353.48\,x} \\[-3pt]
  &\quad +2.388{\times}10^{-6}\,e^{-441.85\,x}
\end{aligned}\!\right)}, \quad x = 1{-}h
\end{equation*}

\noindent\textbf{R3+11} 2.372~ppb\nopagebreak
\begin{equation*}
P = \frac{\pi(a{+}b)\,R_3(h)}{1 - h^{3.5774}\!\left(\!\begin{aligned}
  &\; +1.180{\times}10^{-6}\,e^{-6.3895\,x} \; -5.901{\times}10^{-6}\,e^{-9.9393\,x} \\[-3pt]
  &\quad -5.724{\times}10^{-6}\,e^{-14.909\,x} \; -5.879{\times}10^{-6}\,e^{-17.394\,x} \\[-3pt]
  &\quad +7.871{\times}10^{-7}\,e^{-22.363\,x} \; +9.588{\times}10^{-6}\,e^{-48.305\,x} \\[-3pt]
  &\quad +9.346{\times}10^{-6}\,e^{-96.610\,x} \; -5.901{\times}10^{-6}\,e^{-112.71\,x} \\[-3pt]
  &\quad +4.090{\times}10^{-6}\,e^{-223.63\,x} \; -3.846{\times}10^{-6}\,e^{-357.81\,x} \\[-3pt]
  &\quad +2.345{\times}10^{-6}\,e^{-447.27\,x}
\end{aligned}\!\right)}, \quad x = 1{-}h
\end{equation*}

\noindent\textbf{R3+12} 2.326~ppb\nopagebreak
\begin{equation*}
P = \frac{\pi(a{+}b)\,R_3(h)}{1 - h^{2.7386}\!\left(\!\begin{aligned}
  &\; +1.017{\times}10^{-6}\,e^{-6.5886\,x} \; -5.011{\times}10^{-6}\,e^{-10.249\,x} \\[-3pt]
  &\quad -5.074{\times}10^{-6}\,e^{-15.373\,x} \; -4.830{\times}10^{-6}\,e^{-17.936\,x} \\[-3pt]
  &\quad -4.863{\times}10^{-6}\,e^{-23.060\,x} \; +4.967{\times}10^{-6}\,e^{-32.592\,x} \\[-3pt]
  &\quad +7.825{\times}10^{-6}\,e^{-49.810\,x} \; +8.642{\times}10^{-6}\,e^{-99.619\,x} \\[-3pt]
  &\quad -4.830{\times}10^{-6}\,e^{-116.22\,x} \; +3.482{\times}10^{-6}\,e^{-230.60\,x} \\[-3pt]
  &\quad -3.319{\times}10^{-6}\,e^{-368.96\,x} \; +2.080{\times}10^{-6}\,e^{-461.20\,x}
\end{aligned}\!\right)}, \quad x = 1{-}h
\end{equation*}

\noindent\textbf{R3+13} 1.835~ppb\nopagebreak
\begin{equation*}
P = \frac{\pi(a{+}b)\,R_3(h)}{1 - h^{3.5796}\!\left(\!\begin{aligned}
  &\; +1.111{\times}10^{-6}\,e^{-6.1889\,x} \; -5.274{\times}10^{-6}\,e^{-9.6272\,x} \\[-3pt]
  &\quad -5.274{\times}10^{-6}\,e^{-14.441\,x} \; -5.275{\times}10^{-6}\,e^{-16.848\,x} \\[-3pt]
  &\quad -1.568{\times}10^{-6}\,e^{-21.661\,x} \; -1.402{\times}10^{-7}\,e^{-30.615\,x} \\[-3pt]
  &\quad +1.106{\times}10^{-5}\,e^{-46.788\,x} \; +1.974{\times}10^{-6}\,e^{-93.577\,x} \\[-3pt]
  &\quad +2.321{\times}10^{-6}\,e^{-109.17\,x} \; +4.405{\times}10^{-7}\,e^{-216.61\,x} \\[-3pt]
  &\quad +3.000{\times}10^{-6}\,e^{-346.58\,x} \; -5.290{\times}10^{-6}\,e^{-433.23\,x} \\[-3pt]
  &\quad +3.005{\times}10^{-6}\,e^{-505.43\,x}
\end{aligned}\!\right)}, \quad x = 1{-}h
\end{equation*}

\noindent\textbf{R3+14} 1.495~ppb\nopagebreak
\begin{equation*}
P = \frac{\pi(a{+}b)\,R_3(h)}{1 - h^{3.5796}\!\left(\!\begin{aligned}
  &\; +1.029{\times}10^{-6}\,e^{-6.2073\,x} \; -5.130{\times}10^{-6}\,e^{-9.6558\,x} \\[-3pt]
  &\quad -4.889{\times}10^{-6}\,e^{-14.484\,x} \; -5.043{\times}10^{-6}\,e^{-16.898\,x} \\[-3pt]
  &\quad -4.890{\times}10^{-6}\,e^{-21.725\,x} \; +5.303{\times}10^{-6}\,e^{-30.705\,x} \\[-3pt]
  &\quad +6.894{\times}10^{-6}\,e^{-46.927\,x} \; +5.281{\times}10^{-6}\,e^{-93.854\,x} \\[-3pt]
  &\quad +1.186{\times}10^{-6}\,e^{-109.50\,x} \; -4.891{\times}10^{-6}\,e^{-217.25\,x} \\[-3pt]
  &\quad +9.320{\times}10^{-6}\,e^{-281.95\,x} \; -2.581{\times}10^{-6}\,e^{-347.61\,x} \\[-3pt]
  &\quad -4.980{\times}10^{-6}\,e^{-434.51\,x} \; +3.475{\times}10^{-6}\,e^{-506.93\,x}
\end{aligned}\!\right)}, \quad x = 1{-}h
\end{equation*}

\noindent\textbf{R3+15} \textbf{1.427~ppb}\nopagebreak
\begin{equation*}
P = \frac{\pi(a{+}b)\,R_3(h)}{1 - h^{3.4164}\!\left(\!\begin{aligned}
  &\; +1.069{\times}10^{-6}\,e^{-6.2414\,x} \; -5.345{\times}10^{-6}\,e^{-9.7089\,x} \\[-3pt]
  &\quad -4.028{\times}10^{-6}\,e^{-14.563\,x} \; -5.070{\times}10^{-6}\,e^{-16.991\,x} \\[-3pt]
  &\quad -5.344{\times}10^{-6}\,e^{-21.845\,x} \; +3.808{\times}10^{-6}\,e^{-30.874\,x} \\[-3pt]
  &\quad +9.540{\times}10^{-6}\,e^{-47.185\,x} \; -3.477{\times}10^{-6}\,e^{-94.370\,x} \\[-3pt]
  &\quad +1.034{\times}10^{-5}\,e^{-110.10\,x} \; -3.937{\times}10^{-6}\,e^{-189.32\,x} \\[-3pt]
  &\quad -3.353{\times}10^{-6}\,e^{-218.45\,x} \; +1.059{\times}10^{-5}\,e^{-283.50\,x} \\[-3pt]
  &\quad -3.064{\times}10^{-6}\,e^{-349.52\,x} \; -5.329{\times}10^{-6}\,e^{-436.90\,x} \\[-3pt]
  &\quad +3.678{\times}10^{-6}\,e^{-509.72\,x}
\end{aligned}\!\right)}, \quad x = 1{-}h
\end{equation*}

\medskip\noindent\textbf{R4 family} ($R_4$ base, matching $B(h)$ through $h^{14}$, error budget $T_4 = -1.84 \times 10^{-10}$)\nopagebreak

\noindent In every R4 formula below, $R_4(h)$ stands for $R_3(h)$ (defined above) plus a second Pad\'e layer:
\[
R_4(h) = R_3(h) \;+\; h^{10}\!\cdot\!\frac{-1.979{\times}10^{-7} + 1.864{\times}10^{-9}\,h + 1.238{\times}10^{-7}\,h^2 + 1.083{\times}10^{-7}\,h^3}
       {1 - 2.680\,h + 2.014\,h^2}
\]

\noindent\textbf{R4+7} 3.634~ppb\nopagebreak
\begin{equation*}
P = \frac{\pi(a{+}b)\,R_4(h)}{1 - h^{10.8247}\!\left(\!\begin{aligned}
  &\; -4.354{\times}10^{-6}\,e^{-5.8430\,x} \; -1.389{\times}10^{-5}\,e^{-7.5125\,x} \; +5.850{\times}10^{-6}\,e^{-17.529\,x} \\[-3pt]
  &\quad +9.510{\times}10^{-6}\,e^{-56.794\,x} \; +2.367{\times}10^{-6}\,e^{-132.52\,x} \; +1.474{\times}10^{-7}\,e^{-262.94\,x} \\[-3pt]
  &\quad +3.708{\times}10^{-7}\,e^{-525.87\,x}
\end{aligned}\!\right)}, \quad x = 1{-}h
\end{equation*}

\noindent\textbf{R4+8} 3.612~ppb\nopagebreak
\begin{equation*}
P = \frac{\pi(a{+}b)\,R_4(h)}{1 - h^{11.0642}\!\left(\!\begin{aligned}
  &\; -3.559{\times}10^{-6}\,e^{-5.5058\,x} \; -1.448{\times}10^{-5}\,e^{-7.0789\,x} \; +5.228{\times}10^{-6}\,e^{-16.517\,x} \\[-3pt]
  &\quad +9.424{\times}10^{-6}\,e^{-53.516\,x} \; +2.881{\times}10^{-6}\,e^{-124.87\,x} \; +9.443{\times}10^{-8}\,e^{-247.76\,x} \\[-3pt]
  &\quad +3.257{\times}10^{-7}\,e^{-495.52\,x} \; +8.161{\times}10^{-8}\,e^{-578.11\,x}
\end{aligned}\!\right)}, \quad x = 1{-}h
\end{equation*}

\noindent\textbf{R4+9} 2.258~ppb\nopagebreak
\begin{equation*}
P = \frac{\pi(a{+}b)\,R_4(h)}{1 - h^{13.3435}\!\left(\!\begin{aligned}
  &\; -9.724{\times}10^{-6}\,e^{-4.4909\,x} \; -2.508{\times}10^{-5}\,e^{-5.7740\,x} \; +4.743{\times}10^{-5}\,e^{-10.105\,x} \\[-3pt]
  &\quad -2.873{\times}10^{-5}\,e^{-13.473\,x} \; +1.255{\times}10^{-5}\,e^{-43.652\,x} \; +1.690{\times}10^{-6}\,e^{-101.85\,x} \\[-3pt]
  &\quad +2.215{\times}10^{-6}\,e^{-202.09\,x} \; -2.410{\times}10^{-6}\,e^{-404.18\,x} \; +2.055{\times}10^{-6}\,e^{-471.55\,x}
\end{aligned}\!\right)}, \quad x = 1{-}h
\end{equation*}

\noindent\textbf{R4+10} 1.965~ppb\nopagebreak
\begin{equation*}
P = \frac{\pi(a{+}b)\,R_4(h)}{1 - h^{13.2997}\!\left(\!\begin{aligned}
  &\; -9.174{\times}10^{-6}\,e^{-4.5482\,x} \; -2.497{\times}10^{-5}\,e^{-5.8477\,x} \; +4.358{\times}10^{-5}\,e^{-10.234\,x} \\[-3pt]
  &\quad -2.485{\times}10^{-5}\,e^{-13.645\,x} \; +1.108{\times}10^{-5}\,e^{-44.209\,x} \; +8.037{\times}10^{-6}\,e^{-103.15\,x} \\[-3pt]
  &\quad -6.873{\times}10^{-6}\,e^{-122.80\,x} \; +3.865{\times}10^{-6}\,e^{-204.67\,x} \; -3.208{\times}10^{-6}\,e^{-409.34\,x} \\[-3pt]
  &\quad +2.503{\times}10^{-6}\,e^{-477.56\,x}
\end{aligned}\!\right)}, \quad x = 1{-}h
\end{equation*}

\noindent\textbf{R4+11} 1.539~ppb\nopagebreak
\begin{equation*}
P = \frac{\pi(a{+}b)\,R_4(h)}{1 - h^{2.0246}\!\left(\!\begin{aligned}
  &\; -1.209{\times}10^{-6}\,e^{-5.0779\,x} \; +2.959{\times}10^{-6}\,e^{-6.5287\,x} \\[-3pt]
  &\quad -2.882{\times}10^{-6}\,e^{-11.425\,x} \; -1.421{\times}10^{-5}\,e^{-15.234\,x} \\[-3pt]
  &\quad +9.389{\times}10^{-6}\,e^{-49.357\,x} \; +9.206{\times}10^{-6}\,e^{-115.17\,x} \\[-3pt]
  &\quad +5.084{\times}10^{-6}\,e^{-137.10\,x} \; -1.275{\times}10^{-5}\,e^{-150.29\,x} \\[-3pt]
  &\quad +5.218{\times}10^{-6}\,e^{-228.50\,x} \; -3.140{\times}10^{-6}\,e^{-457.01\,x} \\[-3pt]
  &\quad +2.334{\times}10^{-6}\,e^{-533.18\,x}
\end{aligned}\!\right)}, \quad x = 1{-}h
\end{equation*}

\noindent\textbf{R4+12} 1.136~ppb\nopagebreak
\begin{equation*}
P = \frac{\pi(a{+}b)\,R_4(h)}{1 - h^{0.8519}\!\left(\!\begin{aligned}
  &\; -9.848{\times}10^{-7}\,e^{-6.0323\,x} \; +2.761{\times}10^{-6}\,e^{-7.7558\,x} \\[-3pt]
  &\quad -6.892{\times}10^{-6}\,e^{-13.573\,x} \; -1.121{\times}10^{-5}\,e^{-18.097\,x} \\[-3pt]
  &\quad +6.365{\times}10^{-6}\,e^{-42.226\,x} \; +3.462{\times}10^{-6}\,e^{-58.634\,x} \\[-3pt]
  &\quad +3.073{\times}10^{-6}\,e^{-83.877\,x} \; +7.361{\times}10^{-6}\,e^{-136.81\,x} \\[-3pt]
  &\quad -6.178{\times}10^{-6}\,e^{-162.87\,x} \; +2.682{\times}10^{-6}\,e^{-271.45\,x} \\[-3pt]
  &\quad -1.948{\times}10^{-6}\,e^{-542.90\,x} \; +1.505{\times}10^{-6}\,e^{-633.39\,x}
\end{aligned}\!\right)}, \quad x = 1{-}h
\end{equation*}

\noindent\textbf{R4+13} 1.073~ppb\nopagebreak
\begin{equation*}
P = \frac{\pi(a{+}b)\,R_4(h)}{1 - h^{0.8803}\!\left(\!\begin{aligned}
  &\; -1.061{\times}10^{-6}\,e^{-6.2622\,x} \; +3.105{\times}10^{-6}\,e^{-8.0514\,x} \\[-3pt]
  &\quad -9.878{\times}10^{-6}\,e^{-14.090\,x} \; -8.387{\times}10^{-6}\,e^{-18.787\,x} \\[-3pt]
  &\quad +6.958{\times}10^{-6}\,e^{-43.835\,x} \; +3.713{\times}10^{-6}\,e^{-60.869\,x} \\[-3pt]
  &\quad -3.628{\times}10^{-6}\,e^{-87.074\,x} \; +7.235{\times}10^{-6}\,e^{-93.933\,x} \\[-3pt]
  &\quad +3.915{\times}10^{-6}\,e^{-142.03\,x} \; -3.867{\times}10^{-6}\,e^{-169.08\,x} \\[-3pt]
  &\quad +2.276{\times}10^{-6}\,e^{-281.80\,x} \; -1.737{\times}10^{-6}\,e^{-563.60\,x} \\[-3pt]
  &\quad +1.357{\times}10^{-6}\,e^{-657.53\,x}
\end{aligned}\!\right)}, \quad x = 1{-}h
\end{equation*}

\noindent\textbf{R4+14} 0.936~ppb\nopagebreak
\begin{equation*}
P = \frac{\pi(a{+}b)\,R_4(h)}{1 - h^{1.1845}\!\left(\!\begin{aligned}
  &\; -1.200{\times}10^{-6}\,e^{-6.1369\,x} \; +3.450{\times}10^{-6}\,e^{-7.8902\,x} \\[-3pt]
  &\quad -1.057{\times}10^{-5}\,e^{-13.808\,x} \; -7.698{\times}10^{-6}\,e^{-18.411\,x} \\[-3pt]
  &\quad +5.824{\times}10^{-6}\,e^{-42.958\,x} \; +4.336{\times}10^{-6}\,e^{-59.650\,x} \\[-3pt]
  &\quad +3.679{\times}10^{-6}\,e^{-85.331\,x} \; -7.430{\times}10^{-7}\,e^{-92.053\,x} \\[-3pt]
  &\quad +1.765{\times}10^{-6}\,e^{-139.18\,x} \; +1.599{\times}10^{-6}\,e^{-165.70\,x} \\[-3pt]
  &\quad -8.321{\times}10^{-6}\,e^{-276.16\,x} \; +9.117{\times}10^{-6}\,e^{-312.98\,x} \\[-3pt]
  &\quad -3.454{\times}10^{-6}\,e^{-552.32\,x} \; +2.219{\times}10^{-6}\,e^{-644.37\,x}
\end{aligned}\!\right)}, \quad x = 1{-}h
\end{equation*}

\noindent\textbf{R4+15} 0.786~ppb\nopagebreak
\begin{equation*}
P = \frac{\pi(a{+}b)\,R_4(h)}{1 - h^{1.4531}\!\left(\!\begin{aligned}
  &\; -1.341{\times}10^{-6}\,e^{-6.1633\,x} \; +3.893{\times}10^{-6}\,e^{-7.9242\,x} \\[-3pt]
  &\quad -1.289{\times}10^{-5}\,e^{-13.867\,x} \; -5.414{\times}10^{-6}\,e^{-18.490\,x} \\[-3pt]
  &\quad +5.405{\times}10^{-6}\,e^{-43.143\,x} \; +3.963{\times}10^{-6}\,e^{-59.907\,x} \\[-3pt]
  &\quad +7.588{\times}10^{-6}\,e^{-85.699\,x} \; -2.905{\times}10^{-6}\,e^{-92.449\,x} \\[-3pt]
  &\quad -6.271{\times}10^{-6}\,e^{-139.78\,x} \; +1.069{\times}10^{-5}\,e^{-166.41\,x} \\[-3pt]
  &\quad -1.300{\times}10^{-5}\,e^{-277.35\,x} \; +1.016{\times}10^{-5}\,e^{-314.33\,x} \\[-3pt]
  &\quad +4.677{\times}10^{-6}\,e^{-462.24\,x} \; -8.371{\times}10^{-6}\,e^{-554.69\,x} \\[-3pt]
  &\quad +3.817{\times}10^{-6}\,e^{-647.14\,x}
\end{aligned}\!\right)}, \quad x = 1{-}h
\end{equation*}

\noindent\textbf{R4+16} 0.741~ppb\nopagebreak
\begin{equation*}
P = \frac{\pi(a{+}b)\,R_4(h)}{1 - h^{1.5901}\!\left(\!\begin{aligned}
  &\; -1.372{\times}10^{-6}\,e^{-6.0374\,x} \; +3.917{\times}10^{-6}\,e^{-7.7624\,x} \\[-3pt]
  &\quad -1.210{\times}10^{-5}\,e^{-13.584\,x} \; -6.306{\times}10^{-6}\,e^{-18.112\,x} \\[-3pt]
  &\quad +5.672{\times}10^{-6}\,e^{-42.262\,x} \; +2.917{\times}10^{-6}\,e^{-58.684\,x} \\[-3pt]
  &\quad +7.734{\times}10^{-6}\,e^{-83.949\,x} \; -1.584{\times}10^{-6}\,e^{-90.561\,x} \\[-3pt]
  &\quad -7.266{\times}10^{-6}\,e^{-136.93\,x} \; +1.064{\times}10^{-5}\,e^{-163.01\,x} \\[-3pt]
  &\quad -9.803{\times}10^{-6}\,e^{-271.68\,x} \; +6.820{\times}10^{-6}\,e^{-307.91\,x} \\[-3pt]
  &\quad +4.884{\times}10^{-6}\,e^{-452.81\,x} \; -6.771{\times}10^{-6}\,e^{-543.37\,x} \\[-3pt]
  &\quad +2.338{\times}10^{-6}\,e^{-633.93\,x} \; +2.764{\times}10^{-7}\,e^{-815.05\,x}
\end{aligned}\!\right)}, \quad x = 1{-}h
\end{equation*}

\noindent\textbf{R4+17} 0.809~ppb\nopagebreak
\begin{equation*}
P = \frac{\pi(a{+}b)\,R_4(h)}{1 - h^{2.3358}\!\left(\!\begin{aligned}
  &\; -1.471{\times}10^{-6}\,e^{-5.5986\,x} \; +4.091{\times}10^{-6}\,e^{-7.1982\,x} \\[-3pt]
  &\quad -1.081{\times}10^{-5}\,e^{-12.597\,x} \; -7.696{\times}10^{-6}\,e^{-16.796\,x} \\[-3pt]
  &\quad +4.686{\times}10^{-6}\,e^{-39.190\,x} \; +2.243{\times}10^{-6}\,e^{-50.387\,x} \\[-3pt]
  &\quad -3.188{\times}10^{-7}\,e^{-54.418\,x} \; +1.031{\times}10^{-5}\,e^{-77.847\,x} \\[-3pt]
  &\quad -1.669{\times}10^{-6}\,e^{-83.979\,x} \; -8.981{\times}10^{-6}\,e^{-126.98\,x} \\[-3pt]
  &\quad +1.136{\times}10^{-5}\,e^{-151.16\,x} \; -4.717{\times}10^{-6}\,e^{-251.94\,x} \\[-3pt]
  &\quad +5.001{\times}10^{-7}\,e^{-285.53\,x} \; +8.608{\times}10^{-6}\,e^{-419.89\,x} \\[-3pt]
  &\quad -9.109{\times}10^{-6}\,e^{-503.87\,x} \; +2.522{\times}10^{-6}\,e^{-587.85\,x} \\[-3pt]
  &\quad +4.492{\times}10^{-7}\,e^{-755.81\,x}
\end{aligned}\!\right)}, \quad x = 1{-}h
\end{equation*}

\noindent\textbf{R4+18} 0.685~ppb\nopagebreak
\begin{equation*}
P = \frac{\pi(a{+}b)\,R_4(h)}{1 - h^{3.3367}\!\left(\!\begin{aligned}
  &\; -1.542{\times}10^{-6}\,e^{-5.5843\,x} \; +4.594{\times}10^{-6}\,e^{-7.1798\,x} \\[-3pt]
  &\quad -1.469{\times}10^{-5}\,e^{-12.565\,x} \; -1.110{\times}10^{-5}\,e^{-16.753\,x} \\[-3pt]
  &\quad +9.178{\times}10^{-6}\,e^{-19.545\,x} \; -3.572{\times}10^{-6}\,e^{-39.090\,x} \\[-3pt]
  &\quad +9.883{\times}10^{-6}\,e^{-50.259\,x} \; +1.463{\times}10^{-6}\,e^{-54.279\,x} \\[-3pt]
  &\quad +2.151{\times}10^{-6}\,e^{-77.648\,x} \; +2.618{\times}10^{-6}\,e^{-83.764\,x} \\[-3pt]
  &\quad -8.783{\times}10^{-6}\,e^{-126.65\,x} \; +1.147{\times}10^{-5}\,e^{-150.78\,x} \\[-3pt]
  &\quad -1.930{\times}10^{-6}\,e^{-251.29\,x} \; -2.860{\times}10^{-6}\,e^{-284.80\,x} \\[-3pt]
  &\quad +7.540{\times}10^{-6}\,e^{-418.82\,x} \; -4.332{\times}10^{-6}\,e^{-502.59\,x} \\[-3pt]
  &\quad -1.133{\times}10^{-6}\,e^{-586.35\,x} \; +1.043{\times}10^{-6}\,e^{-753.88\,x}
\end{aligned}\!\right)}, \quad x = 1{-}h
\end{equation*}

\noindent\textbf{R4+19} 0.525~ppb\nopagebreak
\begin{equation*}
P = \frac{\pi(a{+}b)\,R_4(h)}{1 - h^{2.0604}\!\left(\!\begin{aligned}
  &\; -1.324{\times}10^{-6}\,e^{-6.0306\,x} \; +3.832{\times}10^{-6}\,e^{-7.7536\,x} \\[-3pt]
  &\quad -1.175{\times}10^{-5}\,e^{-13.569\,x} \; -1.465{\times}10^{-5}\,e^{-18.092\,x} \\[-3pt]
  &\quad +1.004{\times}10^{-5}\,e^{-21.107\,x} \; -1.662{\times}10^{-6}\,e^{-42.214\,x} \\[-3pt]
  &\quad +8.143{\times}10^{-6}\,e^{-54.275\,x} \; +2.946{\times}10^{-6}\,e^{-58.617\,x} \\[-3pt]
  &\quad -2.490{\times}10^{-6}\,e^{-83.854\,x} \; +5.923{\times}10^{-6}\,e^{-90.459\,x} \\[-3pt]
  &\quad -5.525{\times}10^{-6}\,e^{-136.77\,x} \; +7.013{\times}10^{-6}\,e^{-162.83\,x} \\[-3pt]
  &\quad +1.086{\times}10^{-6}\,e^{-271.38\,x} \; -3.777{\times}10^{-6}\,e^{-307.56\,x} \\[-3pt]
  &\quad +4.737{\times}10^{-6}\,e^{-452.29\,x} \; -3.289{\times}10^{-6}\,e^{-542.75\,x} \\[-3pt]
  &\quad +7.622{\times}10^{-7}\,e^{-633.21\,x} \; -1.405{\times}10^{-7}\,e^{-814.13\,x} \\[-3pt]
  &\quad +1.216{\times}10^{-7}\,e^{-1266.4\,x}
\end{aligned}\!\right)}, \quad x = 1{-}h
\end{equation*}

\noindent\textbf{R4+20} \textbf{0.492~ppb}\nopagebreak
\begin{equation*}
P = \frac{\pi(a{+}b)\,R_4(h)}{1 - h^{2.5199}\!\left(\!\begin{aligned}
  &\; -1.594{\times}10^{-6}\,e^{-5.2674\,x} \; +4.005{\times}10^{-6}\,e^{-6.7724\,x} \\[-3pt]
  &\quad -6.078{\times}10^{-6}\,e^{-11.852\,x} \; -1.789{\times}10^{-5}\,e^{-15.802\,x} \\[-3pt]
  &\quad +7.005{\times}10^{-6}\,e^{-18.436\,x} \; -1.534{\times}10^{-6}\,e^{-36.872\,x} \\[-3pt]
  &\quad +8.604{\times}10^{-6}\,e^{-47.407\,x} \; -5.919{\times}10^{-7}\,e^{-51.199\,x} \\[-3pt]
  &\quad +4.016{\times}10^{-6}\,e^{-73.242\,x} \; +2.309{\times}10^{-6}\,e^{-79.011\,x} \\[-3pt]
  &\quad -4.071{\times}10^{-6}\,e^{-119.47\,x} \; +6.387{\times}10^{-6}\,e^{-142.22\,x} \\[-3pt]
  &\quad -4.451{\times}10^{-7}\,e^{-189.63\,x} \; -7.202{\times}10^{-6}\,e^{-237.03\,x} \\[-3pt]
  &\quad +9.244{\times}10^{-6}\,e^{-268.64\,x} \; -3.986{\times}10^{-6}\,e^{-395.06\,x} \\[-3pt]
  &\quad -1.640{\times}10^{-6}\,e^{-474.07\,x} \; +5.453{\times}10^{-6}\,e^{-553.08\,x} \\[-3pt]
  &\quad -2.320{\times}10^{-6}\,e^{-711.10\,x} \; +3.261{\times}10^{-7}\,e^{-1106.2\,x}
\end{aligned}\!\right)}, \quad x = 1{-}h
\end{equation*}

\medskip\noindent\textbf{R5 family} ($R_5$ base, matching $B(h)$ through $h^{19}$, error budget $T_5 = +1.06 \times 10^{-10}$)\nopagebreak

\noindent In every R5 formula below, $R_5(h)$ stands for $R_4(h)$ (defined above) plus a third Pad\'e layer:
\[
R_5(h) = R_4(h) \;+\; h^{15}\!\cdot\!\frac{-1.472{\times}10^{-7} - 1.167{\times}10^{-7}\,h + 5.780{\times}10^{-8}\,h^2 + 2.056{\times}10^{-7}\,h^3}
       {1 - 3.886\,h + 4.420\,h^2}
\]

\noindent\textbf{R5+8} 3.728~ppb\nopagebreak
\begin{equation*}
P = \frac{\pi(a{+}b)\,R_5(h)}{1 - h^{3.3124}\!\left(\!\begin{aligned}
  &\; +2.001{\times}10^{-6}\,e^{-8.7561\,x} \; -1.549{\times}10^{-5}\,e^{-13.621\,x} \; -1.600{\times}10^{-5}\,e^{-20.431\,x} \\[-3pt]
  &\quad +1.749{\times}10^{-5}\,e^{-23.836\,x} \; +9.425{\times}10^{-6}\,e^{-66.196\,x} \; +2.289{\times}10^{-6}\,e^{-154.46\,x} \\[-3pt]
  &\quad -2.308{\times}10^{-7}\,e^{-306.47\,x} \; +5.228{\times}10^{-7}\,e^{-490.34\,x}
\end{aligned}\!\right)}, \quad x = 1{-}h
\end{equation*}

\noindent\textbf{R5+9} 1.523~ppb\nopagebreak
\begin{equation*}
P = \frac{\pi(a{+}b)\,R_5(h)}{1 - h^{2.9717}\!\left(\!\begin{aligned}
  &\; +1.968{\times}10^{-6}\,e^{-8.8355\,x} \; -1.495{\times}10^{-5}\,e^{-13.744\,x} \; -1.502{\times}10^{-5}\,e^{-20.616\,x} \\[-3pt]
  &\quad +1.567{\times}10^{-5}\,e^{-24.052\,x} \; +1.023{\times}10^{-5}\,e^{-66.797\,x} \; +1.055{\times}10^{-6}\,e^{-155.86\,x} \\[-3pt]
  &\quad +2.044{\times}10^{-6}\,e^{-309.24\,x} \; -1.075{\times}10^{-5}\,e^{-494.79\,x} \; +9.762{\times}10^{-6}\,e^{-515.41\,x}
\end{aligned}\!\right)}, \quad x = 1{-}h
\end{equation*}

\noindent\textbf{R5+10} 1.451~ppb\nopagebreak
\begin{equation*}
P = \frac{\pi(a{+}b)\,R_5(h)}{1 - h^{2.6310}\!\left(\!\begin{aligned}
  &\; +1.862{\times}10^{-6}\,e^{-9.0571\,x} \; -1.473{\times}10^{-5}\,e^{-14.089\,x} \; -1.439{\times}10^{-5}\,e^{-21.133\,x} \\[-3pt]
  &\quad +1.398{\times}10^{-5}\,e^{-24.656\,x} \; +1.324{\times}10^{-6}\,e^{-31.700\,x} \; +9.925{\times}10^{-6}\,e^{-68.472\,x} \\[-3pt]
  &\quad +1.030{\times}10^{-6}\,e^{-159.77\,x} \; +1.935{\times}10^{-6}\,e^{-317.00\,x} \; -1.020{\times}10^{-5}\,e^{-507.20\,x} \\[-3pt]
  &\quad +9.260{\times}10^{-6}\,e^{-528.33\,x}
\end{aligned}\!\right)}, \quad x = 1{-}h
\end{equation*}

\noindent\textbf{R5+11} 1.221~ppb\nopagebreak
\begin{equation*}
P = \frac{\pi(a{+}b)\,R_5(h)}{1 - h^{3.3831}\!\left(\!\begin{aligned}
  &\; +2.088{\times}10^{-6}\,e^{-8.4496\,x} \; -1.445{\times}10^{-5}\,e^{-13.144\,x} \\[-3pt]
  &\quad -1.600{\times}10^{-5}\,e^{-19.716\,x} \; +1.612{\times}10^{-5}\,e^{-23.002\,x} \\[-3pt]
  &\quad -5.104{\times}10^{-7}\,e^{-29.574\,x} \; +1.028{\times}10^{-5}\,e^{-63.879\,x} \\[-3pt]
  &\quad +1.708{\times}10^{-6}\,e^{-149.05\,x} \; +5.312{\times}10^{-7}\,e^{-295.74\,x} \\[-3pt]
  &\quad +1.329{\times}10^{-5}\,e^{-473.18\,x} \; -1.587{\times}10^{-5}\,e^{-492.89\,x} \\[-3pt]
  &\quad +2.824{\times}10^{-6}\,e^{-591.47\,x}
\end{aligned}\!\right)}, \quad x = 1{-}h
\end{equation*}

\noindent\textbf{R5+12} 1.189~ppb\nopagebreak
\begin{equation*}
P = \frac{\pi(a{+}b)\,R_5(h)}{1 - h^{3.1473}\!\left(\!\begin{aligned}
  &\; +2.069{\times}10^{-6}\,e^{-8.8213\,x} \; -1.593{\times}10^{-5}\,e^{-13.722\,x} \\[-3pt]
  &\quad -1.444{\times}10^{-5}\,e^{-20.583\,x} \; +1.745{\times}10^{-5}\,e^{-24.014\,x} \\[-3pt]
  &\quad -1.916{\times}10^{-6}\,e^{-30.875\,x} \; +1.162{\times}10^{-5}\,e^{-66.689\,x} \\[-3pt]
  &\quad -2.060{\times}10^{-6}\,e^{-102.92\,x} \; +2.795{\times}10^{-6}\,e^{-155.61\,x} \\[-3pt]
  &\quad -7.581{\times}10^{-8}\,e^{-308.75\,x} \; +1.365{\times}10^{-5}\,e^{-494.00\,x} \\[-3pt]
  &\quad -1.573{\times}10^{-5}\,e^{-514.58\,x} \; +2.567{\times}10^{-6}\,e^{-617.49\,x}
\end{aligned}\!\right)}, \quad x = 1{-}h
\end{equation*}

\noindent\textbf{R5+13} 0.854~ppb\nopagebreak
\begin{equation*}
P = \frac{\pi(a{+}b)\,R_5(h)}{1 - h^{4.1502}\!\left(\!\begin{aligned}
  &\; +2.514{\times}10^{-6}\,e^{-8.2603\,x} \; -1.758{\times}10^{-5}\,e^{-12.849\,x} \\[-3pt]
  &\quad -1.797{\times}10^{-5}\,e^{-19.274\,x} \; +2.854{\times}10^{-5}\,e^{-22.486\,x} \\[-3pt]
  &\quad -1.034{\times}10^{-5}\,e^{-28.911\,x} \; +1.619{\times}10^{-5}\,e^{-62.448\,x} \\[-3pt]
  &\quad -8.493{\times}10^{-6}\,e^{-96.371\,x} \; +9.155{\times}10^{-6}\,e^{-145.71\,x} \\[-3pt]
  &\quad -1.124{\times}10^{-5}\,e^{-289.11\,x} \; +2.084{\times}10^{-5}\,e^{-385.48\,x} \\[-3pt]
  &\quad +1.272{\times}10^{-6}\,e^{-462.58\,x} \; -1.905{\times}10^{-5}\,e^{-481.85\,x} \\[-3pt]
  &\quad +6.166{\times}10^{-6}\,e^{-578.22\,x}
\end{aligned}\!\right)}, \quad x = 1{-}h
\end{equation*}

\noindent\textbf{R5+14} 0.812~ppb\nopagebreak
\begin{equation*}
P = \frac{\pi(a{+}b)\,R_5(h)}{1 - h^{4.0554}\!\left(\!\begin{aligned}
  &\; +2.587{\times}10^{-6}\,e^{-8.5399\,x} \; -1.988{\times}10^{-5}\,e^{-13.284\,x} \\[-3pt]
  &\quad -1.230{\times}10^{-5}\,e^{-19.926\,x} \; +2.560{\times}10^{-5}\,e^{-23.247\,x} \\[-3pt]
  &\quad -1.112{\times}10^{-5}\,e^{-29.890\,x} \; +1.894{\times}10^{-5}\,e^{-64.562\,x} \\[-3pt]
  &\quad -1.639{\times}10^{-5}\,e^{-99.632\,x} \; +2.778{\times}10^{-5}\,e^{-139.48\,x} \\[-3pt]
  &\quad -1.431{\times}10^{-5}\,e^{-150.64\,x} \; -8.047{\times}10^{-6}\,e^{-298.90\,x} \\[-3pt]
  &\quad +1.629{\times}10^{-5}\,e^{-398.53\,x} \; +4.675{\times}10^{-6}\,e^{-478.23\,x} \\[-3pt]
  &\quad -1.943{\times}10^{-5}\,e^{-498.16\,x} \; +5.600{\times}10^{-6}\,e^{-597.79\,x}
\end{aligned}\!\right)}, \quad x = 1{-}h
\end{equation*}

\noindent\textbf{R5+15} 0.821~ppb\nopagebreak
\begin{equation*}
P = \frac{\pi(a{+}b)\,R_5(h)}{1 - h^{4.1407}\!\left(\!\begin{aligned}
  &\; +2.547{\times}10^{-6}\,e^{-8.3605\,x} \; -1.846{\times}10^{-5}\,e^{-13.005\,x} \\[-3pt]
  &\quad -1.669{\times}10^{-5}\,e^{-19.508\,x} \; +2.894{\times}10^{-5}\,e^{-22.759\,x} \\[-3pt]
  &\quad -1.136{\times}10^{-5}\,e^{-29.262\,x} \; +1.722{\times}10^{-5}\,e^{-63.206\,x} \\[-3pt]
  &\quad -1.096{\times}10^{-5}\,e^{-97.540\,x} \; +1.453{\times}10^{-5}\,e^{-136.56\,x} \\[-3pt]
  &\quad -7.378{\times}10^{-6}\,e^{-147.48\,x} \; +5.407{\times}10^{-6}\,e^{-195.08\,x} \\[-3pt]
  &\quad -1.493{\times}10^{-5}\,e^{-292.62\,x} \; +2.437{\times}10^{-5}\,e^{-390.16\,x} \\[-3pt]
  &\quad -5.148{\times}10^{-7}\,e^{-468.19\,x} \; -1.915{\times}10^{-5}\,e^{-487.70\,x} \\[-3pt]
  &\quad +6.436{\times}10^{-6}\,e^{-585.24\,x}
\end{aligned}\!\right)}, \quad x = 1{-}h
\end{equation*}

\noindent\textbf{R5+16} 0.599~ppb\nopagebreak
\begin{equation*}
P = \frac{\pi(a{+}b)\,R_5(h)}{1 - h^{6.5194}\!\left(\!\begin{aligned}
  &\; +3.860{\times}10^{-6}\,e^{-6.1489\,x} \; -1.547{\times}10^{-5}\,e^{-9.5649\,x} \\[-3pt]
  &\quad -3.005{\times}10^{-5}\,e^{-14.347\,x} \; +4.072{\times}10^{-5}\,e^{-16.739\,x} \\[-3pt]
  &\quad -1.482{\times}10^{-5}\,e^{-21.521\,x} \; +7.743{\times}10^{-6}\,e^{-46.486\,x} \\[-3pt]
  &\quad +1.019{\times}10^{-5}\,e^{-71.737\,x} \; -1.606{\times}10^{-6}\,e^{-100.43\,x} \\[-3pt]
  &\quad -1.528{\times}10^{-5}\,e^{-108.47\,x} \; +2.405{\times}10^{-5}\,e^{-143.47\,x} \\[-3pt]
  &\quad -2.349{\times}10^{-5}\,e^{-215.21\,x} \; +2.685{\times}10^{-5}\,e^{-286.95\,x} \\[-3pt]
  &\quad -7.344{\times}10^{-6}\,e^{-344.34\,x} \; -9.678{\times}10^{-6}\,e^{-358.69\,x} \\[-3pt]
  &\quad +4.265{\times}10^{-6}\,e^{-430.42\,x} \; +4.337{\times}10^{-8}\,e^{-1434.7\,x}
\end{aligned}\!\right)}, \quad x = 1{-}h
\end{equation*}

\noindent\textbf{R5+17} 0.649~ppb\nopagebreak
\begin{equation*}
P = \frac{\pi(a{+}b)\,R_5(h)}{1 - h^{5.8454}\!\left(\!\begin{aligned}
  &\; +3.440{\times}10^{-6}\,e^{-6.8177\,x} \; -1.677{\times}10^{-5}\,e^{-10.605\,x} \\[-3pt]
  &\quad -2.638{\times}10^{-5}\,e^{-15.908\,x} \; +3.921{\times}10^{-5}\,e^{-18.559\,x} \\[-3pt]
  &\quad -1.558{\times}10^{-5}\,e^{-23.862\,x} \; +1.280{\times}10^{-5}\,e^{-51.542\,x} \\[-3pt]
  &\quad -1.311{\times}10^{-6}\,e^{-79.540\,x} \; +2.072{\times}10^{-5}\,e^{-111.36\,x} \\[-3pt]
  &\quad -2.611{\times}10^{-5}\,e^{-120.26\,x} \; +1.730{\times}10^{-5}\,e^{-159.08\,x} \\[-3pt]
  &\quad -1.022{\times}10^{-5}\,e^{-190.90\,x} \; +6.087{\times}10^{-6}\,e^{-238.62\,x} \\[-3pt]
  &\quad -1.025{\times}10^{-5}\,e^{-318.16\,x} \; -2.077{\times}10^{-5}\,e^{-381.79\,x} \\[-3pt]
  &\quad +3.767{\times}10^{-5}\,e^{-397.70\,x} \; -1.052{\times}10^{-5}\,e^{-477.24\,x} \\[-3pt]
  &\quad +6.927{\times}10^{-7}\,e^{-795.40\,x}
\end{aligned}\!\right)}, \quad x = 1{-}h
\end{equation*}

\noindent\textbf{R5+18} 0.573~ppb\nopagebreak
\begin{equation*}
P = \frac{\pi(a{+}b)\,R_5(h)}{1 - h^{6.2305}\!\left(\!\begin{aligned}
  &\; +3.872{\times}10^{-6}\,e^{-6.6360\,x} \; -1.844{\times}10^{-5}\,e^{-10.323\,x} \\[-3pt]
  &\quad -2.657{\times}10^{-5}\,e^{-15.484\,x} \; +4.331{\times}10^{-5}\,e^{-18.065\,x} \\[-3pt]
  &\quad -1.881{\times}10^{-5}\,e^{-23.226\,x} \; +1.316{\times}10^{-5}\,e^{-50.168\,x} \\[-3pt]
  &\quad -3.421{\times}10^{-6}\,e^{-71.767\,x} \; +5.715{\times}10^{-6}\,e^{-77.420\,x} \\[-3pt]
  &\quad +1.811{\times}10^{-6}\,e^{-108.39\,x} \; -1.161{\times}10^{-5}\,e^{-117.06\,x} \\[-3pt]
  &\quad +1.929{\times}10^{-5}\,e^{-154.84\,x} \; -2.178{\times}10^{-5}\,e^{-232.26\,x} \\[-3pt]
  &\quad +2.408{\times}10^{-5}\,e^{-309.68\,x} \; +7.919{\times}10^{-6}\,e^{-371.61\,x} \\[-3pt]
  &\quad -2.293{\times}10^{-5}\,e^{-387.10\,x} \; -1.289{\times}10^{-6}\,e^{-402.58\,x} \\[-3pt]
  &\quad +5.660{\times}10^{-6}\,e^{-464.52\,x} \; +3.327{\times}10^{-8}\,e^{-1548.4\,x}
\end{aligned}\!\right)}, \quad x = 1{-}h
\end{equation*}

\noindent\textbf{R5+19} 0.571~ppb\nopagebreak
\begin{equation*}
P = \frac{\pi(a{+}b)\,R_5(h)}{1 - h^{6.2123}\!\left(\!\begin{aligned}
  &\; +3.918{\times}10^{-6}\,e^{-6.6892\,x} \; -1.910{\times}10^{-5}\,e^{-10.405\,x} \\[-3pt]
  &\quad -2.503{\times}10^{-5}\,e^{-15.608\,x} \; +4.269{\times}10^{-5}\,e^{-18.210\,x} \\[-3pt]
  &\quad -1.938{\times}10^{-5}\,e^{-23.412\,x} \; +1.537{\times}10^{-5}\,e^{-50.571\,x} \\[-3pt]
  &\quad -1.763{\times}10^{-5}\,e^{-72.343\,x} \; +2.011{\times}10^{-5}\,e^{-78.041\,x} \\[-3pt]
  &\quad -1.073{\times}10^{-5}\,e^{-109.26\,x} \; +1.427{\times}10^{-6}\,e^{-118.00\,x} \\[-3pt]
  &\quad +2.569{\times}10^{-5}\,e^{-156.08\,x} \; -1.601{\times}10^{-5}\,e^{-171.69\,x} \\[-3pt]
  &\quad -7.689{\times}10^{-6}\,e^{-234.12\,x} \; +1.031{\times}10^{-5}\,e^{-312.16\,x} \\[-3pt]
  &\quad +1.625{\times}10^{-5}\,e^{-374.60\,x} \; -2.418{\times}10^{-5}\,e^{-390.20\,x} \\[-3pt]
  &\quad -2.234{\times}10^{-7}\,e^{-405.81\,x} \; +4.180{\times}10^{-6}\,e^{-468.25\,x} \\[-3pt]
  &\quad +3.372{\times}10^{-8}\,e^{-1560.8\,x}
\end{aligned}\!\right)}, \quad x = 1{-}h
\end{equation*}

\noindent\textbf{R5+20} \textbf{0.508~ppb}\nopagebreak
\begin{equation*}
P = \frac{\pi(a{+}b)\,R_5(h)}{1 - h^{6.6769}\!\left(\!\begin{aligned}
  &\; +4.826{\times}10^{-6}\,e^{-6.5273\,x} \; -2.400{\times}10^{-5}\,e^{-10.154\,x} \\[-3pt]
  &\quad -1.398{\times}10^{-5}\,e^{-15.230\,x} \; +3.611{\times}10^{-5}\,e^{-17.769\,x} \\[-3pt]
  &\quad -2.063{\times}10^{-5}\,e^{-22.845\,x} \; +1.836{\times}10^{-5}\,e^{-49.346\,x} \\[-3pt]
  &\quad -3.036{\times}10^{-5}\,e^{-70.591\,x} \; +2.987{\times}10^{-5}\,e^{-76.152\,x} \\[-3pt]
  &\quad +6.103{\times}10^{-6}\,e^{-106.61\,x} \; -2.172{\times}10^{-5}\,e^{-115.14\,x} \\[-3pt]
  &\quad +4.819{\times}10^{-5}\,e^{-152.30\,x} \; -3.003{\times}10^{-5}\,e^{-167.53\,x} \\[-3pt]
  &\quad -1.871{\times}10^{-5}\,e^{-228.45\,x} \; +3.540{\times}10^{-5}\,e^{-304.61\,x} \\[-3pt]
  &\quad +2.202{\times}10^{-5}\,e^{-365.53\,x} \; -3.665{\times}10^{-5}\,e^{-380.76\,x} \\[-3pt]
  &\quad -2.688{\times}10^{-5}\,e^{-395.99\,x} \; +2.303{\times}10^{-5}\,e^{-456.91\,x} \\[-3pt]
  &\quad -1.018{\times}10^{-6}\,e^{-761.52\,x} \; +9.109{\times}10^{-8}\,e^{-1523.0\,x}
\end{aligned}\!\right)}, \quad x = 1{-}h
\end{equation*}

\medskip\noindent\textbf{R6 family} ($R_6$ base, matching $B(h)$ through $h^{24}$, error budget $T_6 \approx 1.7 \times 10^{-11}$)\nopagebreak

\noindent In every R6 formula below, $R_6(h)$ stands for $R_5(h)$ (defined above) plus a fourth and final Pad\'e layer:
\[
R_6(h) = R_5(h) \;+\; h^{20}\!\cdot\!\frac{-1.057{\times}10^{-6} - 1.828{\times}10^{-6}\,h - 4.965{\times}10^{-7}\,h^2 + 3.382{\times}10^{-6}\,h^3}
       {1 - 5.589\,h + 9.324\,h^2}
\]
The full chain is: $R_6(h) = 1 + 3h/(10{+}\sqrt{4{-}3h}) + h^5\,\psi_3(h) + h^{10}\,\psi_4(h) + h^{15}\,\psi_5(h) + h^{20}\,\psi_6(h)$, where each $\psi_k$ is the cubic/quadratic fraction shown in the corresponding family header above.

\noindent\textbf{R6+7} 4.071~ppb\nopagebreak
\begin{equation*}
P = \frac{\pi(a{+}b)\,R_6(h)}{1 - h^{1.3359}\!\left(\!\begin{aligned}
  &\; -1.412{\times}10^{-6}\,e^{-5.1725\,x} \; +2.869{\times}10^{-6}\,e^{-6.6503\,x} \; -1.693{\times}10^{-5}\,e^{-15.517\,x} \\[-3pt]
  &\quad +1.019{\times}10^{-5}\,e^{-50.276\,x} \; +4.912{\times}10^{-6}\,e^{-117.31\,x} \; -1.850{\times}10^{-7}\,e^{-232.76\,x} \\[-3pt]
  &\quad +5.639{\times}10^{-7}\,e^{-465.52\,x}
\end{aligned}\!\right)}, \quad x = 1{-}h
\end{equation*}

\noindent\textbf{R6+8} 3.613~ppb\nopagebreak
\begin{equation*}
P = \frac{\pi(a{+}b)\,R_6(h)}{1 - h^{2.7310}\!\left(\!\begin{aligned}
  &\; -1.458{\times}10^{-6}\,e^{-5.7547\,x} \; +4.788{\times}10^{-6}\,e^{-7.3989\,x} \; -1.738{\times}10^{-5}\,e^{-12.948\,x} \\[-3pt]
  &\quad -4.965{\times}10^{-7}\,e^{-17.264\,x} \; +1.074{\times}10^{-5}\,e^{-55.936\,x} \; +3.345{\times}10^{-6}\,e^{-130.52\,x} \\[-3pt]
  &\quad +4.581{\times}10^{-8}\,e^{-258.96\,x} \; +4.174{\times}10^{-7}\,e^{-517.92\,x}
\end{aligned}\!\right)}, \quad x = 1{-}h
\end{equation*}

\noindent\textbf{R6+9} 1.649~ppb\nopagebreak
\begin{equation*}
P = \frac{\pi(a{+}b)\,R_6(h)}{1 - h^{2.7423}\!\left(\!\begin{aligned}
  &\; -1.425{\times}10^{-6}\,e^{-5.8121\,x} \; +4.524{\times}10^{-6}\,e^{-7.4727\,x} \; -1.696{\times}10^{-5}\,e^{-13.077\,x} \\[-3pt]
  &\quad -8.135{\times}10^{-7}\,e^{-17.436\,x} \; +1.129{\times}10^{-5}\,e^{-56.493\,x} \; +2.462{\times}10^{-6}\,e^{-131.82\,x} \\[-3pt]
  &\quad +9.718{\times}10^{-7}\,e^{-261.54\,x} \; -1.020{\times}10^{-6}\,e^{-523.09\,x} \; +9.766{\times}10^{-7}\,e^{-610.27\,x}
\end{aligned}\!\right)}, \quad x = 1{-}h
\end{equation*}

\noindent\textbf{R6+10} 1.332~ppb\nopagebreak
\begin{equation*}
P = \frac{\pi(a{+}b)\,R_6(h)}{1 - h^{2.4825}\!\left(\!\begin{aligned}
  &\; -1.589{\times}10^{-6}\,e^{-5.5515\,x} \; +4.498{\times}10^{-6}\,e^{-7.1377\,x} \; -1.226{\times}10^{-5}\,e^{-12.491\,x} \\[-3pt]
  &\quad -5.702{\times}10^{-6}\,e^{-16.655\,x} \; +1.099{\times}10^{-5}\,e^{-53.961\,x} \; +3.682{\times}10^{-6}\,e^{-125.91\,x} \\[-3pt]
  &\quad -4.290{\times}10^{-6}\,e^{-249.82\,x} \; +5.505{\times}10^{-6}\,e^{-283.13\,x} \; -2.812{\times}10^{-6}\,e^{-499.64\,x} \\[-3pt]
  &\quad +1.984{\times}10^{-6}\,e^{-582.91\,x}
\end{aligned}\!\right)}, \quad x = 1{-}h
\end{equation*}

\noindent\textbf{R6+11} 0.955~ppb\nopagebreak
\begin{equation*}
P = \frac{\pi(a{+}b)\,R_6(h)}{1 - h^{1.6227}\!\left(\!\begin{aligned}
  &\; -1.659{\times}10^{-6}\,e^{-6.2671\,x} \; +4.777{\times}10^{-6}\,e^{-8.0577\,x} \\[-3pt]
  &\quad -1.656{\times}10^{-5}\,e^{-14.101\,x} \; -1.437{\times}10^{-6}\,e^{-18.801\,x} \\[-3pt]
  &\quad +1.261{\times}10^{-5}\,e^{-56.404\,x} \; -1.508{\times}10^{-6}\,e^{-60.916\,x} \\[-3pt]
  &\quad +3.987{\times}10^{-6}\,e^{-142.14\,x} \; -6.721{\times}10^{-6}\,e^{-282.02\,x} \\[-3pt]
  &\quad +7.531{\times}10^{-6}\,e^{-319.62\,x} \; -2.959{\times}10^{-6}\,e^{-564.04\,x} \\[-3pt]
  &\quad +1.940{\times}10^{-6}\,e^{-658.04\,x}
\end{aligned}\!\right)}, \quad x = 1{-}h
\end{equation*}

\noindent\textbf{R6+12} 0.830~ppb\nopagebreak
\begin{equation*}
P = \frac{\pi(a{+}b)\,R_6(h)}{1 - h^{0.8837}\!\left(\!\begin{aligned}
  &\; -9.690{\times}10^{-7}\,e^{-6.5416\,x} \; +2.961{\times}10^{-6}\,e^{-8.4106\,x} \\[-3pt]
  &\quad -1.104{\times}10^{-5}\,e^{-14.719\,x} \; -9.051{\times}10^{-6}\,e^{-19.625\,x} \\[-3pt]
  &\quad +4.899{\times}10^{-6}\,e^{-29.437\,x} \; +3.787{\times}10^{-6}\,e^{-58.874\,x} \\[-3pt]
  &\quad +6.102{\times}10^{-6}\,e^{-63.584\,x} \; +3.459{\times}10^{-6}\,e^{-148.36\,x} \\[-3pt]
  &\quad -6.487{\times}10^{-6}\,e^{-294.37\,x} \; +7.364{\times}10^{-6}\,e^{-333.62\,x} \\[-3pt]
  &\quad -2.929{\times}10^{-6}\,e^{-588.74\,x} \; +1.898{\times}10^{-6}\,e^{-686.86\,x}
\end{aligned}\!\right)}, \quad x = 1{-}h
\end{equation*}

\noindent\textbf{R6+13} 0.623~ppb\nopagebreak
\begin{equation*}
P = \frac{\pi(a{+}b)\,R_6(h)}{1 - h^{1.6030}\!\left(\!\begin{aligned}
  &\; -1.302{\times}10^{-6}\,e^{-6.2193\,x} \; +3.826{\times}10^{-6}\,e^{-7.9963\,x} \\[-3pt]
  &\quad -1.306{\times}10^{-5}\,e^{-13.994\,x} \; -6.283{\times}10^{-6}\,e^{-18.658\,x} \\[-3pt]
  &\quad +2.874{\times}10^{-6}\,e^{-27.987\,x} \; +4.675{\times}10^{-6}\,e^{-55.974\,x} \\[-3pt]
  &\quad +5.688{\times}10^{-6}\,e^{-60.452\,x} \; +3.542{\times}10^{-6}\,e^{-141.05\,x} \\[-3pt]
  &\quad -4.388{\times}10^{-6}\,e^{-279.87\,x} \; +4.971{\times}10^{-6}\,e^{-317.19\,x} \\[-3pt]
  &\quad -1.683{\times}10^{-6}\,e^{-559.74\,x} \; +1.128{\times}10^{-6}\,e^{-653.03\,x} \\[-3pt]
  &\quad +1.144{\times}10^{-8}\,e^{-2239.0\,x}
\end{aligned}\!\right)}, \quad x = 1{-}h
\end{equation*}

\noindent\textbf{R6+14} 0.700~ppb\nopagebreak
\begin{equation*}
P = \frac{\pi(a{+}b)\,R_6(h)}{1 - h^{1.2615}\!\left(\!\begin{aligned}
  &\; -1.126{\times}10^{-6}\,e^{-6.4095\,x} \; +3.395{\times}10^{-6}\,e^{-8.2408\,x} \\[-3pt]
  &\quad -1.245{\times}10^{-5}\,e^{-14.421\,x} \; -7.250{\times}10^{-6}\,e^{-19.228\,x} \\[-3pt]
  &\quad +3.898{\times}10^{-6}\,e^{-28.843\,x} \; +4.499{\times}10^{-6}\,e^{-57.685\,x} \\[-3pt]
  &\quad +5.605{\times}10^{-6}\,e^{-62.300\,x} \; +3.566{\times}10^{-6}\,e^{-145.37\,x} \\[-3pt]
  &\quad -6.960{\times}10^{-6}\,e^{-288.43\,x} \; +8.273{\times}10^{-6}\,e^{-326.88\,x} \\[-3pt]
  &\quad -2.643{\times}10^{-6}\,e^{-499.94\,x} \; +6.588{\times}10^{-7}\,e^{-576.85\,x} \\[-3pt]
  &\quad +4.997{\times}10^{-7}\,e^{-673.00\,x} \; +3.101{\times}10^{-8}\,e^{-1346.0\,x}
\end{aligned}\!\right)}, \quad x = 1{-}h
\end{equation*}

\noindent\textbf{R6+15} 0.477~ppb\nopagebreak
\begin{equation*}
P = \frac{\pi(a{+}b)\,R_6(h)}{1 - h^{1.9386}\!\left(\!\begin{aligned}
  &\; -1.384{\times}10^{-6}\,e^{-6.1721\,x} \; +4.062{\times}10^{-6}\,e^{-7.9356\,x} \\[-3pt]
  &\quad -1.392{\times}10^{-5}\,e^{-13.887\,x} \; -6.884{\times}10^{-6}\,e^{-18.516\,x} \\[-3pt]
  &\quad +1.023{\times}10^{-5}\,e^{-27.775\,x} \; -7.382{\times}10^{-6}\,e^{-32.404\,x} \\[-3pt]
  &\quad +7.741{\times}10^{-6}\,e^{-55.549\,x} \; +4.207{\times}10^{-6}\,e^{-59.993\,x} \\[-3pt]
  &\quad +2.698{\times}10^{-6}\,e^{-139.98\,x} \; +8.206{\times}10^{-7}\,e^{-277.75\,x} \\[-3pt]
  &\quad -2.021{\times}10^{-7}\,e^{-314.78\,x} \; -3.481{\times}10^{-6}\,e^{-481.43\,x} \\[-3pt]
  &\quad +6.160{\times}10^{-6}\,e^{-555.49\,x} \; -2.799{\times}10^{-6}\,e^{-648.07\,x} \\[-3pt]
  &\quad +1.345{\times}10^{-7}\,e^{-1296.1\,x}
\end{aligned}\!\right)}, \quad x = 1{-}h
\end{equation*}

\noindent\textbf{R6+16} \textbf{0.018~ppb} \textbf{--- (this work, superseded by SCOR)}\nopagebreak
\begin{equation*}
P = \frac{\pi(a{+}b)\,R_6(h)}{1 - h^{7.6874}\!\left(\!\begin{aligned}
  &\; -7.477{\times}10^{-5}\,e^{-6.4814\,x} \; +8.480{\times}10^{-4}\,e^{-8.8090\,x} \\[-3pt]
  &\quad -1.375{\times}10^{-3}\,e^{-10.180\,x} \; +2.151{\times}10^{-3}\,e^{-14.569\,x} \\[-3pt]
  &\quad -1.958{\times}10^{-3}\,e^{-16.657\,x} \; +5.984{\times}10^{-4}\,e^{-22.580\,x} \\[-3pt]
  &\quad -2.023{\times}10^{-4}\,e^{-26.260\,x} \; +8.698{\times}10^{-6}\,e^{-57.325\,x} \\[-3pt]
  &\quad +2.452{\times}10^{-6}\,e^{-91.868\,x} \; +1.741{\times}10^{-6}\,e^{-200.09\,x} \\[-3pt]
  &\quad -5.234{\times}10^{-6}\,e^{-372.92\,x} \; +5.453{\times}10^{-6}\,e^{-395.77\,x} \\[-3pt]
  &\quad -8.113{\times}10^{-7}\,e^{-694.96\,x} \; +5.347{\times}10^{-7}\,e^{-832.67\,x} \\[-3pt]
  &\quad -2.118{\times}10^{-8}\,e^{-2816.7\,x} \; +2.201{\times}10^{-8}\,e^{-3499.8\,x}
\end{aligned}\!\right)}, \quad x = 1{-}h
\end{equation*}

\noindent\textbf{R6+17} 0.521~ppb\nopagebreak
\begin{equation*}
P = \frac{\pi(a{+}b)\,R_6(h)}{1 - h^{1.9074}\!\left(\!\begin{aligned}
  &\; -1.461{\times}10^{-6}\,e^{-6.3141\,x} \; +4.393{\times}10^{-6}\,e^{-8.1181\,x} \\[-3pt]
  &\quad -1.707{\times}10^{-5}\,e^{-14.207\,x} \; -1.709{\times}10^{-6}\,e^{-18.942\,x} \\[-3pt]
  &\quad +1.978{\times}10^{-6}\,e^{-28.414\,x} \; +2.953{\times}10^{-6}\,e^{-56.827\,x} \\[-3pt]
  &\quad +8.046{\times}10^{-6}\,e^{-61.373\,x} \; -4.872{\times}10^{-6}\,e^{-122.75\,x} \\[-3pt]
  &\quad +1.099{\times}10^{-5}\,e^{-143.20\,x} \; -1.305{\times}10^{-5}\,e^{-208.37\,x} \\[-3pt]
  &\quad +4.123{\times}10^{-6}\,e^{-246.25\,x} \; +9.425{\times}10^{-6}\,e^{-263.39\,x} \\[-3pt]
  &\quad +7.343{\times}10^{-6}\,e^{-284.14\,x} \; -1.353{\times}10^{-5}\,e^{-322.02\,x} \\[-3pt]
  &\quad +8.415{\times}10^{-6}\,e^{-568.27\,x} \; -7.552{\times}10^{-6}\,e^{-662.98\,x} \\[-3pt]
  &\quad +1.582{\times}10^{-6}\,e^{-852.41\,x}
\end{aligned}\!\right)}, \quad x = 1{-}h
\end{equation*}

\noindent\textbf{R6+18} 0.358~ppb (coarse grid; true peak 115~ppb at $e > 0.999$)\nopagebreak
\begin{equation*}
P = \frac{\pi(a{+}b)\,R_6(h)}{1 - h^{2.6332}\!\left(\!\begin{aligned}
  &\; -1.889{\times}10^{-6}\,e^{-5.8138\,x} \; +5.261{\times}10^{-6}\,e^{-7.4749\,x} \\[-3pt]
  &\quad -1.695{\times}10^{-5}\,e^{-13.081\,x} \; -1.142{\times}10^{-6}\,e^{-17.441\,x} \\[-3pt]
  &\quad +1.177{\times}10^{-6}\,e^{-26.162\,x} \; +5.943{\times}10^{-7}\,e^{-52.324\,x} \\[-3pt]
  &\quad +9.234{\times}10^{-6}\,e^{-56.510\,x} \; +2.754{\times}10^{-6}\,e^{-131.86\,x} \\[-3pt]
  &\quad -4.673{\times}10^{-6}\,e^{-261.62\,x} \; +1.089{\times}10^{-5}\,e^{-296.50\,x} \\[-3pt]
  &\quad -5.656{\times}10^{-8}\,e^{-313.94\,x} \; -7.711{\times}10^{-6}\,e^{-348.83\,x} \\[-3pt]
  &\quad +2.464{\times}10^{-6}\,e^{-523.24\,x} \; -1.079{\times}10^{-6}\,e^{-610.45\,x} \\[-3pt]
  &\quad +2.868{\times}10^{-6}\,e^{-627.89\,x} \; -2.722{\times}10^{-6}\,e^{-784.86\,x} \\[-3pt]
  &\quad -1.686{\times}10^{-7}\,e^{-872.07\,x} \; +1.150{\times}10^{-6}\,e^{-959.27\,x}
\end{aligned}\!\right)}, \quad x = 1{-}h
\end{equation*}

\noindent\textbf{R6+19} 0.479~ppb\nopagebreak
\begin{equation*}
P = \frac{\pi(a{+}b)\,R_6(h)}{1 - h^{2.2492}\!\left(\!\begin{aligned}
  &\; -1.725{\times}10^{-6}\,e^{-6.2158\,x} \; +5.127{\times}10^{-6}\,e^{-7.9918\,x} \\[-3pt]
  &\quad -1.950{\times}10^{-5}\,e^{-13.986\,x} \; -6.247{\times}10^{-7}\,e^{-18.647\,x} \\[-3pt]
  &\quad +4.278{\times}10^{-6}\,e^{-21.755\,x} \; -2.854{\times}10^{-6}\,e^{-27.971\,x} \\[-3pt]
  &\quad +7.848{\times}10^{-6}\,e^{-55.942\,x} \; +9.104{\times}10^{-6}\,e^{-60.418\,x} \\[-3pt]
  &\quad -7.917{\times}10^{-6}\,e^{-74.590\,x} \; +6.837{\times}10^{-6}\,e^{-120.84\,x} \\[-3pt]
  &\quad +1.165{\times}10^{-6}\,e^{-140.97\,x} \; -1.053{\times}10^{-5}\,e^{-205.12\,x} \\[-3pt]
  &\quad +8.362{\times}10^{-7}\,e^{-242.42\,x} \; +1.068{\times}10^{-5}\,e^{-259.29\,x} \\[-3pt]
  &\quad +1.017{\times}10^{-5}\,e^{-279.71\,x} \; -1.555{\times}10^{-5}\,e^{-317.01\,x} \\[-3pt]
  &\quad +9.247{\times}10^{-6}\,e^{-559.42\,x} \; -8.320{\times}10^{-6}\,e^{-652.66\,x} \\[-3pt]
  &\quad +1.741{\times}10^{-6}\,e^{-839.14\,x}
\end{aligned}\!\right)}, \quad x = 1{-}h
\end{equation*}

\noindent\textbf{R6+20} 0.769~ppb\nopagebreak
\begin{equation*}
P = \frac{\pi(a{+}b)\,R_6(h)}{1 - h^{2.1787}\!\left(\!\begin{aligned}
  &\; -1.912{\times}10^{-6}\,e^{-6.1961\,x} \; +5.652{\times}10^{-6}\,e^{-7.9664\,x} \\[-3pt]
  &\quad -2.116{\times}10^{-5}\,e^{-13.941\,x} \; +2.092{\times}10^{-6}\,e^{-18.588\,x} \\[-3pt]
  &\quad +5.199{\times}10^{-6}\,e^{-21.686\,x} \; -6.652{\times}10^{-6}\,e^{-27.882\,x} \\[-3pt]
  &\quad +1.139{\times}10^{-5}\,e^{-55.765\,x} \; +1.340{\times}10^{-5}\,e^{-60.226\,x} \\[-3pt]
  &\quad -1.817{\times}10^{-5}\,e^{-74.353\,x} \; +1.383{\times}10^{-5}\,e^{-120.45\,x} \\[-3pt]
  &\quad +7.616{\times}10^{-6}\,e^{-140.53\,x} \; -1.820{\times}10^{-5}\,e^{-167.29\,x} \\[-3pt]
  &\quad +2.752{\times}10^{-6}\,e^{-204.47\,x} \; -9.927{\times}10^{-6}\,e^{-241.65\,x} \\[-3pt]
  &\quad +1.423{\times}10^{-5}\,e^{-258.46\,x} \; +1.452{\times}10^{-5}\,e^{-278.82\,x} \\[-3pt]
  &\quad -1.695{\times}10^{-5}\,e^{-316.00\,x} \; +7.662{\times}10^{-6}\,e^{-557.65\,x} \\[-3pt]
  &\quad -6.862{\times}10^{-6}\,e^{-650.59\,x} \; +1.494{\times}10^{-6}\,e^{-836.47\,x}
\end{aligned}\!\right)}, \quad x = 1{-}h
\end{equation*}

\noindent\textbf{R6+21} 0.922~ppb\nopagebreak
\begin{equation*}
P = \frac{\pi(a{+}b)\,R_6(h)}{1 - h^{2.1577}\!\left(\!\begin{aligned}
  &\; -1.538{\times}10^{-6}\,e^{-5.8982\,x} \; +4.461{\times}10^{-6}\,e^{-7.5834\,x} \\[-3pt]
  &\quad -1.422{\times}10^{-5}\,e^{-13.271\,x} \; -4.676{\times}10^{-6}\,e^{-17.695\,x} \\[-3pt]
  &\quad +1.718{\times}10^{-6}\,e^{-20.644\,x} \; -1.337{\times}10^{-6}\,e^{-26.542\,x} \\[-3pt]
  &\quad +9.155{\times}10^{-6}\,e^{-53.084\,x} \; +6.809{\times}10^{-6}\,e^{-57.330\,x} \\[-3pt]
  &\quad -7.501{\times}10^{-6}\,e^{-70.778\,x} \; +5.939{\times}10^{-6}\,e^{-114.66\,x} \\[-3pt]
  &\quad +8.324{\times}10^{-6}\,e^{-133.77\,x} \; -1.020{\times}10^{-5}\,e^{-159.25\,x} \\[-3pt]
  &\quad +4.129{\times}10^{-7}\,e^{-194.64\,x} \; -1.091{\times}10^{-5}\,e^{-230.03\,x} \\[-3pt]
  &\quad +1.230{\times}10^{-5}\,e^{-246.04\,x} \; +1.064{\times}10^{-5}\,e^{-265.42\,x} \\[-3pt]
  &\quad -1.074{\times}10^{-5}\,e^{-300.81\,x} \; +4.333{\times}10^{-7}\,e^{-389.28\,x} \\[-3pt]
  &\quad +3.135{\times}10^{-6}\,e^{-530.84\,x} \; -3.050{\times}10^{-6}\,e^{-619.31\,x} \\[-3pt]
  &\quad +8.394{\times}10^{-7}\,e^{-796.26\,x}
\end{aligned}\!\right)}, \quad x = 1{-}h
\end{equation*}

\noindent\textbf{R6+22} 0.441~ppb\nopagebreak
\begin{equation*}
P = \frac{\pi(a{+}b)\,R_6(h)}{1 - h^{2.2359}\!\left(\!\begin{aligned}
  &\; -1.706{\times}10^{-6}\,e^{-6.2243\,x} \; +5.079{\times}10^{-6}\,e^{-8.0027\,x} \\[-3pt]
  &\quad -1.939{\times}10^{-5}\,e^{-14.005\,x} \; -1.351{\times}10^{-6}\,e^{-18.673\,x} \\[-3pt]
  &\quad +6.091{\times}10^{-6}\,e^{-21.785\,x} \; -6.273{\times}10^{-6}\,e^{-28.010\,x} \\[-3pt]
  &\quad +5.028{\times}10^{-6}\,e^{-37.346\,x} \; -2.570{\times}10^{-6}\,e^{-56.019\,x} \\[-3pt]
  &\quad +1.341{\times}10^{-5}\,e^{-60.501\,x} \; -2.174{\times}10^{-6}\,e^{-74.692\,x} \\[-3pt]
  &\quad -1.754{\times}10^{-6}\,e^{-121.00\,x} \; +1.201{\times}10^{-5}\,e^{-141.17\,x} \\[-3pt]
  &\quad -9.670{\times}10^{-6}\,e^{-168.06\,x} \; -3.664{\times}10^{-6}\,e^{-205.40\,x} \\[-3pt]
  &\quad +7.324{\times}10^{-6}\,e^{-224.08\,x} \; -1.062{\times}10^{-5}\,e^{-242.75\,x} \\[-3pt]
  &\quad +1.143{\times}10^{-5}\,e^{-259.64\,x} \; +1.200{\times}10^{-5}\,e^{-280.10\,x} \\[-3pt]
  &\quad -1.621{\times}10^{-5}\,e^{-317.44\,x} \; +1.049{\times}10^{-5}\,e^{-560.19\,x} \\[-3pt]
  &\quad -9.429{\times}10^{-6}\,e^{-653.56\,x} \; +1.932{\times}10^{-6}\,e^{-840.29\,x}
\end{aligned}\!\right)}, \quad x = 1{-}h
\end{equation*}

\medskip\noindent\textbf{Total: 55 formulas} across 4 tower levels. The complete machine-readable data (all rates, coefficients, ratios, and errors to full precision) is available in \texttt{best\_results.csv} in the code repository.

\end{document}\n